\documentclass[11pt,a4paper]{article}
\usepackage[margin=27mm]{geometry}
\usepackage{amsmath,amssymb,amsthm,mathtools,bm,microtype}
\usepackage[colorlinks=true,linkcolor=blue,citecolor=blue,urlcolor=blue]{hyperref}
\numberwithin{equation}{section}
\newtheorem{theorem}{Theorem}[section]
\newtheorem{proposition}[theorem]{Proposition}
\newtheorem{lemma}[theorem]{Lemma}
\theoremstyle{remark}
\newtheorem{remark}[theorem]{Remark}
\newcommand{\R}{\mathbb R}
\newcommand{\T}{\mathbb T}
\newcommand{\id}{\mathrm{id}}
\newcommand{\norm}[1]{\left\lVert #1\right\rVert}
\newcommand{\abs}[1]{\left\lvert #1\right\rvert}
\newcommand{\ip}[2]{\left\langle #1,#2\right\rangle_0}
\newcommand{\Graph}{\mathsf E}
\newcommand{\PN}{P_{\!N}}
\newcommand{\Xzero}{\mathcal X_0}
\newcommand{\Xone}{\mathcal X_1}
\newcommand{\Azero}{\mathcal A_0}
\newcommand{\Aone}{\mathcal A_1}
\newcommand{\ad}{\operatorname{ad}}
\newcommand{\diag}{\operatorname{diag}}
\title{Linear failure and cubic convergence for surface evolution\\
with relaxed minimal deformation}
\author{Kaiwen Jin\\[2pt]\small\href{mailto:yui.ui.siki@gmail.com}{\texttt{yui.ui.siki@gmail.com}}}
\date{}
\begin{document}
\maketitle
\begin{abstract}
The finite element scheme for relaxed minimal-deformation surface evolution
prescribes normal velocity while allowing tangential mesh motion.
We establish contrasting results for its original exact-integral formulation
with the discontinuous geometric unit normals of the current elements.
For linear elements, one fixed smooth torus, a tangential rigid-rotation
field, and an admissible alternating mesh family give a positive lower
bound for the uniform-in-time \(L^2\) error of the lifted flow map,
on an interval where the discrete solution exists uniquely and remains
nondegenerate. For cubic elements, every admissible mesh family converges
for every smooth exact solution with smooth inverse on a fixed interval:
the original discrete solution exists uniquely throughout that interval,
and its lifted map error is at most \(C_T h^3\sqrt{1+|\log h|}\).
The cubic proof first controls shape error by a signed-distance-weighted
relative volume and then estimates weakly tangent map error.
These results partially solve the low-degree convergence question;
the quadratic case remains unresolved.
\end{abstract}

\section{Introduction}\label{sec:intro}
Tangential redistribution can improve a surface mesh without changing
its normal velocity. For a parametrized method, however, the resulting
flow map is itself part of the approximation. Its convergence requires
more than control of the shape of the surface.
The relaxed minimal-deformation formulation of Gao, Li and
Tang~\cite{GLT2026} uses a harmonic-map relaxation from the initial
surface to the evolving one. Its finite element scheme uses the
geometric normal of the computed surface directly.

For smooth solutions whose flow maps and inverse maps remain smooth,
Theorem~2.1 of the author manuscript of Gao, Li and Tang~\cite{GLT2026}
proves a lifted, uniform-in-time \(L^2\) error bound of order \(h^k\)
for polynomial degree \(k\geq4\). Section~7 of that manuscript leaves
convergence for \(k=1,2,3\) open. The question concerns the same
position-gradient equation, exact fixed-reference integrals, and
raw geometric normals, which can jump across current element edges.
The theorem and section locators for \cite{GLT2026} used here refer
to the linked author manuscript.

We prove failure for degree one and convergence for degree three.
For the negative result, the initial surface is a fixed smooth torus
and the external field is one fixed small rigid rotation tangent to it.
Thus the continuous relaxed minimal-deformation solution is stationary.
An admissible shape-regular, quasi-uniform mesh family alternates two
unequal azimuthal spacings. There are fixed \(T_*,c_*>0\) such that
the original semidiscrete solutions exist uniquely with nondegenerate
faces on \([0,T_*]\), while
\[
 \liminf_{h\to0}\sup_{0\leq t\leq T_*}
           \norm{X_h^\ell(t)-\id}_{L^2(\Gamma^0)}\geq c_* .
\]
The lift is the closest-point lift from the initial nodal surface.
The construction is given in Theorem~\ref{thm:failure}.

For the positive result, fix any smooth exact solution on \([0,T]\)
with smooth inverse, and any admissible cubic nodal mesh family under
the original uniform parametrization bounds. For all sufficiently
small \(h\), the unmodified semidiscrete solution exists uniquely
with nondegenerate elements through \(T\), and
\[
 \sup_{0\leq t\leq T}\norm{X_h^\ell(t)-X(t)}_{L^2(\Gamma^0)}
          \leq C_T h^3\sqrt{1+|\log h|}.
\]
Theorem~\ref{thm:cubic} retains every smooth-solution and mesh-family
hypothesis of the original cubic question. The constant is independent
of \(h\). This paper gives a partial solution of AIM Problem~541:
the quadratic case remains unresolved.

The linear counterexample separates map error from surface shape.
The alternating species create a displacement proportional to the mesh
width whose reference gradient can remain nonzero. A scalar label
of the weak normal constraint records this displacement. At the initial
torus its rate is \(-\varepsilon\cos\theta/6\), where \(\theta\) is
the meridional angle and \(\varepsilon\) the fixed rotation amplitude.
A common analytic interval and a time modulus for the mean velocity
control the label through the fast initial layer. The original rotation
energy then converts its fixed-time gradient effect into a positive
\(L^2\) map-error lower bound.

For cubic elements a single map-error estimate would mix this tangent
drift with the error in surface shape. We instead use two estimates.
A relative volume weighted by signed distance controls shape error.
Its derivative is paired with a projection defect that retains the
current area ratio, because the algorithm's weak constraint uses the
fixed reference measure. The first area variation annihilates pure
normal slopes, leaving a vanishing tangent-to-shape feedback.
This gives a normal-error bound strong enough to close the tangent
energy, with the logarithmic factor from a two-dimensional discrete
Sobolev estimate. The resulting differential closeness permits
continuation of the actual mixed system to the final time.

Section~\ref{sec:mesh} gives the original equations and the linear mesh
construction. Sections~\ref{sec:operators} and \ref{sec:interval}
control its weak operators and common existence interval;
Section~\ref{sec:label} proves linear failure.
Section~\ref{sec:cubic} treats general smooth data with cubic elements.
All physical fields are real. Fourier arguments use their
complexification, with conventions stated beside their use.
All discrete integrals use the fixed initial surface exactly.

\section{The torus and its triangulations}\label{sec:mesh}

We use the position-gradient relaxed minimal-deformation equations
\begin{equation}\label{eq:continuous}
 X_t=v,\qquad
 (v-U_{\rm ext}(X,t))\cdot(n\circ X)=0,\qquad
 v-\Delta_{\Gamma^0}X=\kappa(n\circ X),\qquad X(0)=\id .
\end{equation}
The Laplace--Beltrami operator and all reference integrals are on the
initial surface \(\Gamma^0\). The unit normal \(n\) belongs to the current
surface. Thus \(\kappa\) is a scalar constraint multiplier.

For degree \(k\), let \(S_h\) be the continuous degree-\(k\) parametric
space on a nodal interpolant \(\Gamma_h^0\). The semidiscrete equations are exactly
\begin{align}
 &X_{h,t}=v_h,\qquad X_h(0)=\id,\label{eq:discrete-time}\\
 &\int_{\Gamma_h^0}(v_h-U_{\rm ext}(X_h,t))\cdot N_h\,\chi_h=0
       &&(\chi_h\in S_h),\label{eq:discrete-normal}\\
 &\int_{\Gamma_h^0}v_h\cdot\eta_h+
       \int_{\Gamma_h^0}\nabla_{\Gamma_h^0}X_h:
                         \nabla_{\Gamma_h^0}\eta_h
       =\int_{\Gamma_h^0}\kappa_h N_h\cdot\eta_h
       &&(\eta_h\in S_h^3).\label{eq:discrete-force}
\end{align}
Here \(N_h\) is the pullback of the full geometric unit normal of each
current element. It can jump across edges. All integrals are exact.
In this section and Sections~\ref{sec:operators}--\ref{sec:label},
we take \(k=1\), so the reference elements are planar triangles.
Section~\ref{sec:cubic} uses the same equations for \(k=3\).

\subsection{Data and the mesh family}

Fix \(R=2\), \(r=2/5\), let \(\T=\R/(2\pi\mathbb Z)\), and set
\begin{equation}\label{eq:torus}
 \Psi(\theta,\phi)=
 ((R+r\cos\theta)\cos\phi,(R+r\cos\theta)\sin\phi,r\sin\theta),
 \qquad (\theta,\phi)\in\T^2.
\end{equation}
The orientation is outward. Write \(\rho_\theta=R+r\cos\theta\),
\(\Omega x=(-x_2,x_1,0)\), and take one fixed, sufficiently small
\(\varepsilon>0\) with
\begin{equation}\label{eq:field}
 U_{\rm ext}(x,t)=\varepsilon\Omega x .
\end{equation}
This field is tangent to the torus. Its metric is
\(\diag(r^2,\rho_\theta^2)\), and direct differentiation gives
\[
 \Delta_{\Gamma^0}\Psi
   =-\left(\frac1r+\frac{\cos\theta}{\rho_\theta}\right)n .
\]
Thus \(X=\id\), \(v=0\), and \(\kappa=1/r+\cos\theta/\rho_\theta\)
solve \eqref{eq:continuous} smoothly on every fixed interval.

For \(n\in\mathbb N\), put
\[
 \delta=\frac{2\pi}{4n},\qquad s=4\delta,\qquad
 \theta_i=is,\qquad a=\frac13,\quad b=\frac23 .
\]
The two azimuthal labels in cell \(j\), \(0\leq j<4n\), are
\[
 \phi_{j,+}=j\delta-a\delta/2,\qquad
 \phi_{j,-}=j\delta+a\delta/2 .
\]
Take nodal values of \(\Psi\) at these angles. In two neighboring rows
denote the nodes in one cell by \(A=(i,+)\), \(B=(i,-)\),
\(C=(i+1,+)\), \(D=(i+1,-)\). Primes indicate the next azimuthal cell.
Use the oriented faces
\begin{equation}\label{eq:faces}
 ABD,\qquad ADC,\qquad BA'C',\qquad BC'D .
\end{equation}
The short and long azimuthal widths are \(a\delta\) and \(b\delta\).
Let \(Y_h\) be the identity on the resulting planar surface.

\begin{lemma}\label{lem:admissible}
For sufficiently large \(n\), these are admissible shape-regular,
quasi-uniform degree-one nodal meshes, with \(h\asymp\delta\).
The closest-point projection \(a_h:\Gamma_h^0\to\Gamma^0\) is bijective,
its face Jacobians are bounded above and below uniformly, and
\[
 \norm{Y_h^\ell-\id}_{L^\infty(\Gamma^0)}\leq C\delta^2,
 \qquad Y_h^\ell=Y_h\circ a_h^{-1}.
\]
Here \(K_f\) is the flat triangle with the same vertices as \(K\).
The degree-one element parametrization \(F_K:K_f\to K\) is the identity.
\end{lemma}
\begin{proof}
The limiting edge directions are the nonparallel meridional and
azimuthal tangents, of lengths \(rs\), \(\rho_\theta a\delta\), and
\(\rho_\theta b\delta\). Their normalized lengths and angles stay
in compact nondegenerate sets. The distance and normal errors are
\(O(\delta^2)\) and \(O(\delta)\), by trigonometric expansion.

The closest projection is explicit. Put
\(\rho_{xy}=(x_1^2+x_2^2)^{1/2}\) and
\(d=((\rho_{xy}-R)^2+x_3^2)^{1/2}\). Then
\[
 a(x)=\left(
 \left[R+\frac{r(\rho_{xy}-R)}d\right]\frac{x_1}{\rho_{xy}},
 \left[R+\frac{r(\rho_{xy}-R)}d\right]\frac{x_2}{\rho_{xy}},
 \frac{rx_3}d\right).
\]
Both denominators are uniformly positive near the meshes. Each
azimuthal quadrilateral is planar. At fixed azimuth \(\phi\), its
cylindrical radius is the affine interpolation of its two row radii
multiplied by
\[
 c_\phi=\frac{\cos(d_\phi/2)}{\cos(\phi-\phi_{\rm mid})}
        =1+O(\delta^2),
\]
where \(d_\phi\) is its angular width. The determinant of the
cross-section vector about \((R,0)\) and its meridional chord is
\(r^2\sin s+O(\delta^2s)>0\). Its projected cross-section angle is
strictly increasing, as is the azimuth. The lifted angular coordinates
cover both periods once. This proves bijectivity. The tangent estimates
give the Jacobian bounds and the identity error.
\end{proof}

\begin{theorem}\label{thm:failure}
For the fixed torus \eqref{eq:torus}, one fixed field
\eqref{eq:field}, and the mesh family \eqref{eq:faces}, there are
\(T_*>0\), \(c_*>0\), and \(n_0\) such that the original system
\eqref{eq:discrete-time}--\eqref{eq:discrete-force} has a unique solution
with nondegenerate evolving elements on \([0,T_*]\) for \(n\geq n_0\),
but
\begin{equation}\label{eq:main-failure}
 \liminf_{n\to\infty}\ \sup_{0\leq t\leq T_*}
       \norm{X_h^\ell(t)-\id}_{L^2(\Gamma^0)}\geq c_* .
\end{equation}
Here \(X_h^\ell=X_h\circ a_h^{-1}\) is the closest-point lift and
\(h\asymp\delta=2\pi/(4n)\). All constants and the field are fixed
before the mesh limit.
\end{theorem}

The theorem concerns the parametrized flow map. Its proof needs a
common existence interval: an initial velocity calculation alone
would not exclude disappearance of the effect in an initial layer.
We establish that interval in Section~\ref{sec:interval} and complete
the proof in Section~\ref{sec:label}.

\subsection{The full local equation and cyclic coordinates}

Let \(M\), \(K\), and \(T(X)\) denote the exact vector mass, reference
stiffness, and weak normal matrices. In physical nodal coordinates,
\begin{equation}\label{eq:matrix-original}
 X_t=V,\qquad MV+KX=T(X)^*\kappa,\qquad
 T(X)(V-\varepsilon\Omega X)=0 .
\end{equation}
Initially \(X=Y\), the nodal identity. These matrices depend smoothly
on \(X\) while all scaled edge cross products stay nonzero.

Let \(N_{\rm ref}\) insert the smooth reference nodal normal. On a
small geometric neighborhood, every current face normal has a common
positive component in this direction. Hence \(TN_{\rm ref}\) is a
definite weighted scalar mass plus \(O(\delta)\): replacing a nodal
reference normal by a common reference normal on each face costs
\(O(\delta)\) in its mass norm, and the remaining weights are positive.
Its Hermitian part has a uniform lower bound. This argument holds on
the full mesh, not only on cyclic modes. Thus \(T\) has full row rank
and \(TM^{-1}T^*\) is invertible.
Eliminating \(\kappa\) gives a smooth finite-dimensional ODE. Picard
iteration on a closed coordinate ball gives a local solution: a short
time makes the integral map a contraction, and the same Lipschitz
bound gives uniqueness.

Cyclic rotations preserve \eqref{eq:matrix-original}. Full uniqueness
therefore gives the exact representation
\begin{equation}\label{eq:cyclic}
 X_{i,j,\pm}=\mathcal R_{\phi_{j,\pm}}(y_i\pm\delta w_i),
 \qquad
 y_i(0)=y_i^0=(R+r\cos\theta_i,0,r\sin\theta_i),\quad w_i(0)=0 ,
\end{equation}
where \(\mathcal R_\phi\) is rotation about the third axis.
All three Cartesian components of both profiles remain present.

\section{The geometric constraint and its weak operators}\label{sec:operators}

\subsection{Row norms and exact element formulas}

The row inner product is \(\ip pq=s\sum_i p_i^*q_i\).
Let \(\tau\) be the forward row permutation and set
\[
 D=\frac{\tau-I}{s},\qquad
 D^*=\frac{\tau^{-1}-I}{s}=-\tau^{-1}D,\qquad
 \lambda_k=\frac{e^{iks}-1}{s}.
\]
Stars refer to this uniform row inner product. A varying-mass adjoint
is written explicitly as \(M^{-1}A^*M\).
With \(\widehat q_k=n^{-1}\sum_iq_i e^{-ik\theta_i}\), define
\begin{equation}\label{eq:row-sobolev}
 \norm q_{H^\sigma}^2
   =2\pi\sum_k(1+|\lambda_k|^2)^\sigma|\widehat q_k|^2 .
\end{equation}
Signed modular frequencies are used. For \(|k|\leq n/2\),
\(|\lambda_k|\geq 2|k|/\pi\), and
\(|\lambda_{k+l}|\leq|\lambda_k|+|\lambda_l|\), including aliasing.
Fourier convolution and Cauchy--Schwarz give the uniform \(H^1\)
embedding into \(\ell^\infty\) and one-dimensional product estimates.

For a physical species vector write \(q_\pm=q_b\pm q_o\), with
label rotations understood, and put \(q_c=q_o/\delta\). Its reference
squared piecewise affine \(H^1\) norm is uniformly equivalent to
\[
 \norm{q_b}_{H^1}^2+\norm{q_c}_0^2+
                      \delta^2\norm{Dq_c}_0^2
 \asymp\norm{q_b}_{H^1}^2+\norm{q_c}_0^2 .
\]
Here \(\delta\norm D\leq1/2\). We shall use the bases
\begin{equation}\label{eq:bases}
 \Xzero=H^4\times H^3,\quad \Xone=H^5\times H^4,\qquad
 \Azero=H^3\times H^2,\quad \Aone=H^4\times H^3 .
\end{equation}
The macro component is the mean of the two nodal species; the optical
component is half their difference. The scaled optical coordinate divides
that difference by \(\delta\). Geometry is \(\mathbf e=(y-y^0,w)\),
so its physical error has mean \(y-y^0\) and optical part \(\delta w\).
The velocity spaces have one fewer row derivative than the corresponding
geometry spaces; each scaled optical component has one fewer than its mean.

Fix \(p=4\). For a row space \(B\) and \(\rho>0\), write
\begin{equation}\label{eq:analytic-norm}
 \varpi_j=\frac{(j+1)^p\rho^j}{j!},\qquad
 \norm q_{\rho,p;B}^2=\sum_{j\geq0}\varpi_j^2\norm{D^jq}_B^2 .
\end{equation}
Coefficient norms have the same weights with coefficient row
differences.

There are \(4n\) azimuthal copies. Multiply each assembled one-cell
weak matrix by \(4n/s=(\pi/2)\delta^{-2}\). In the row inner product
this represents exactly the full physical cyclic energy and leaves
the equations unchanged. Thus normalized \(M,T\) are order zero.
On a triangle of reference area \(A\), the scalar mass is \(A/12\)
times the matrix with diagonal \(2\) and off-diagonal \(1\).
The stiffness is \(K=J^*J\), where \(J\) is the exact reference
gradient including area and metric factors.

For scaled current edges \(g_l=(X_l-X_0)/\delta\), put
\(c=g_1\times g_2\), \(n=c/|c|\), and
\(d_l=(h_l-h_0)/\delta\). The exact normal derivative is
\begin{equation}\label{eq:normal-derivative}
 DN[h]=\frac{(I-nn^{\mathsf T})(d_1\times g_2+g_1\times d_2)}{|c|}.
\end{equation}
We keep this full normal map. Its coefficients are analytic on a
compact nondegenerate set of scaled edges. Rotation remainders
retain their powers of \(\delta\) by
\[
 \frac{\mathcal R_{\beta\delta}-I}{\delta}
       =\beta\int_0^1\Omega\mathcal R_{t\beta\delta}\,dt .
\]

\subsection{Row moments and the normal graph}

For \((Aq)_i=\sum_k a_i(k)q_{i+k}\), put
\(m_\ell(A)_i=\sum_k k^\ell a_i(k)\).
All kernels below have exponentially summable distance moments and
bounded coefficient differences at the displayed regularities.
Class \(0\) imposes no zero condition; classes \(1\) and \(2\) mean
\[
 m_0=\delta c_0,\qquad\hbox{and}\qquad
 m_0=\delta^2c_0,\quad m_1=\delta c_1 ,
\]
respectively. The factors remain present after coefficient differences.
Physical first and second orders are \(\delta^{-1}\) times class \(1\)
and \(\delta^{-2}\) times class \(2\).

Exactly, without commuting matrices,
\begin{align}
 m_0(AB)_i&=\sum_k a_i(k)m_0(B)_{i+k},\nonumber\\
 m_1(AB)_i&=\sum_k k a_i(k)m_0(B)_{i+k}
                    +\sum_k a_i(k)m_1(B)_{i+k}.\label{eq:mom-products}
\end{align}
Subtract the value of each right coefficient at \(i\). A difference
supplies \(s=4\delta\). Thus two class \(1\) factors give class \(2\).
For a class \(2\) factor next to class \(0\), expand once more: the
first moment contributes \(\delta s\), and the second-difference
remainder contributes \(s^2\). Class \(2\) is preserved in either order.

The adjoint moments are
\[
 m_0(A^*)_i=\sum_l a_{i-l}(l)^*,\qquad
 m_1(A^*)_i=-\sum_l l\,a_{i-l}(l)^* .
\]
The finite Laurent identity
\(\tau^{-l}=I-l(\tau-I)+(\tau-I)^2R_l(\tau)\), with
\(\norm{R_l}_{\ell^1}\leq C(1+|l|)^2\), yields
\begin{equation}\label{eq:adjoint-moment}
 m_0(A^*)=m_0(A)^*-sD_{\rm coeff}(m_1(A)^*)
       +s^2\sum_lR_l(\tau)D_{\rm coeff}^2(a(l)^*) .
\end{equation}
Its minus sign corresponds to our forward difference. The analogous
first-moment identity proves adjoint stability of the classes.
Periodic moments are read on lifted periodized kernels; wrapping
tails are exponential and smaller than each fixed required power of
\(\delta\).

A uniformly invertible localized order-zero matrix has a localized
inverse. Conjugate by
\(W_j=\exp(\gamma\,\operatorname{dist}(i,j))\). Weighted kernel
moments give \(\norm{WAW^{-1}-A}\leq C\gamma\), so a fixed small
\(\gamma\) preserves invertibility by a Neumann series. One-row inputs
then give exponential kernel decay for \(A^{-1}\). If
\(A^+=\tau A\tau^{-1}\) and \(\delta_c A=(A^+-A)/s\), exactly
\begin{equation}\label{eq:inverse-difference}
 \delta_c(A^{-1})=-(A^+)^{-1}(\delta_c A)A^{-1}.
\end{equation}
This also supplies the coefficient differences of the global inverse.

The actual mass-normal projection is
\begin{equation}\label{eq:normal-projector}
 \PN=M^{-1}T^*(TM^{-1}T^*)^{-1}T,\qquad Q=I-\PN .
\end{equation}
The full normal-column gap supplies its inverses. Freezing at output
row \(i\) changes each kernel by
\(C\delta(1+\operatorname{dist})\) times an exponential. At zero
row phase the coherent frozen off-diagonal blocks vanish, since the
alternating affine function on a whole rectangle satisfies
\(\int\chi=0\), \(\int\chi^2=A_{\rm rect}/3\).
A zero-sum cross kernel acts on neighboring differences and supplies
the factor \(\delta\) required by the scaled optical norm.
More explicitly, in a coherent frozen strip let \(\zeta\) be the row
Fourier phase and set \(d_\delta(\zeta)=(1+|\zeta-1|^2/s^2)^{1/2}\).
The full physical state consists of both three-component species values,
with mean and optical parts as above. The alternating affine function
\(\chi\) has nodal values \(+1\) and \(-1\) on the two species.
At \(\zeta=1\), its rectangle integrals make both the mass and weak-normal
matrices block diagonal in mean/optical coordinates. The Schur inverse
therefore has the same zero cross blocks, and smooth dependence gives
\(|P_{bo}(\zeta)|+|P_{ob}(\zeta)|\leq C|\zeta-1|\).
Conjugating by \(\diag(d_\delta I_3,\delta^{-1}I_3)\) is bounded, since
\[
 \frac{|\zeta-1|}{\delta d_\delta}\leq C,\qquad
       \delta d_\delta|\zeta-1|\leq C.
\]
The \(O(\delta)\) actual-kernel freezing remainder is bounded before
dividing the optical coordinate by \(\delta\). This gives a bounded
order-zero remainder in the same mixed norm.
The identity \([D,A^{-1}]=-A^{-1}[D,A]A^{-1}\) supplies differentiated
bounds. In particular \(\PN,Q\) are uniformly bounded in macro
\(H^1\) plus scaled-optical \(L^2\), and in their differentiated versions
when the corresponding coefficient norms are retained.

Use the fixed aligned reference frame
\[
 e_\theta=(-\sin\theta,0,\cos\theta),\quad
 e_\phi=(0,1,0),\quad n_0=(\cos\theta,0,\sin\theta).
\]
Let \(C_{\rm all}\) insert all four tangent coordinates
\(\mathbf a=(u,\omega)\), with physical optical size \(\delta\omega\).
Let \(L_{\rm ref}\) extract them and let \(N_{\rm ref}\) insert the
reference normal coordinates. Set
\begin{equation}\label{eq:graph}
 R_N=N_{\rm ref}(TN_{\rm ref})^{-1},\qquad
 \Graph=C_{\rm all}-R_NTC_{\rm all}.
\end{equation}
Then \(T\Graph=0\), \(L_{\rm ref}\Graph=I\), and
\(\Graph'[h]=-R_NT'[h]\Graph\).
Every weak-tangent vector equals \(\Graph L_{\rm ref}\) of itself.
The global graph and its inverse have uniform mixed bounds. In scaled
state coordinates their type is
\begin{equation}\label{eq:state-type}
 \begin{pmatrix}\text{order }0&\delta^2\text{ order }1\\
                 \text{order }1&\text{order }0\end{pmatrix}.
\end{equation}
Consequently
\begin{equation}\label{eq:graph-high}
 \norm{\Graph\mathbf a}_{\Xzero}
 \leq C_{\rm low}\bigl(\norm{\mathbf a}_{\Aone}
            +\norm{\mathbf e}_{\Xone}\norm{\mathbf a}_{\Azero}\bigr).
\end{equation}
Here \(\Xzero,\Xone\) are the geometry bases \((H^4,H^3)\),
\((H^5,H^4)\), whereas \(\Azero,\Aone\) are the velocity bases
\((H^3,H^2)\), \((H^4,H^3)\). Low constants depend on fixed reference
bounds and the stated current low norms, never on \(n\).
The high geometry factor is written explicitly.

\subsection{The curved cancellation and the initial velocity}

For a scalar multiplier define \(H_\kappa h=T'[h]^*\kappa\).
The alternating curvature scale is \(\kappa_o=\lambda/\delta\).
After removing the common factor \(\delta^{-2}\), the physical
mean/optical blocks of \(H_{\lambda/\delta}\) satisfy
\begin{equation}\label{eq:curved-moments}
 \begin{array}{c|ccc}
       &m_0&m_1&\text{kernel size}\\ \hline
 bb&O(\delta^2)&O(\delta)&O(1)\\
 bo,\ ob&O(\delta)&O(1)&O(1)\\
 oo&O(1)&O(1)&O(1).
 \end{array}
\end{equation}
Here the alternating signs are included in the multiplier, and all
entries are factored identities with bounded smooth coefficients.

To prove them, freeze the lower-row divided meridional chord
\(p=(y_{i+1}-y_i)/\delta\), \(q=\Omega y_i\), and \(w_i\).
The leading short and long rectangles have chords
\(p,aq-2w_i\) and \(p,bq+2w_i\). They are planar, even when
the two directions are not orthogonal. The alternating P1 function
\(\chi\) is affine azimuthally. A coarse row test \(L_\nu\) is affine
meridionally, and exact integration on the two triangles gives
\begin{equation}\label{eq:rectangle-zero}
                  \int_{\rm rectangle}\chi L_\nu=0 .
\end{equation}
A constant aligned macro input changes the scaled edges first at
order \(\delta\), through label rotations. That first variation is
affine azimuthally, so its normal variation is constant on each
leading rectangle. Equation~\eqref{eq:rectangle-zero} removes its
order-\(\delta\) pairing. Reference weights, neighboring multipliers,
test rotations and normal coefficients then change it by another
\(\delta\). This proves the \(bb\) zero sum.
For the \(bb\) first moment the input is affine in row level; its
leading normal variation is again constant on the whole rectangle.
One curvature or coefficient difference remains. For \(bo\), a
constant optical input is affine azimuthally and the same integral
vanishes. For \(ob\), the constant leading macro input is a
translation, so its normal variation is zero. The \(oo\) bound follows
from \eqref{eq:normal-derivative}.

These arguments factor the remainders by exact divided differences
and the rotation integral. For example, the \(bb\) zero-sum
coefficient uses macro first and second differences, optical first
differences and multiplier first differences. A coefficient row
difference raises each corresponding count by one and retains both
mesh factors. Thus the conclusions apply to the actual curved
triangles and their coefficient derivatives.

We need a distinct cancellation for the moving constraint.
For \(W_\pm=W_b\pm\delta W_c\), put
\(B_W h=M_{\rm scalar}^{-1}T'[h]W\). After removing
\(\delta^{-1}\), its physical scalar-output/vector-input blocks have
\begin{equation}\label{eq:connection-moments}
 \begin{array}{c|cc}
       &m_0&m_1\\ \hline
 bb&O(\delta)&O(1)\\
 bo&O(1)&O(1)\\
 ob&O(\delta^2)&O(\delta)\\
 oo&O(\delta)&O(1).
 \end{array}
\end{equation}
The optical test at just one meridional row is important. On a short
rectangle, a constant vector has optical test integral \(A_{\rm tri}/3\)
at the lower row and \(-A_{\rm tri}/3\) at the upper row; on a long
rectangle the signs reverse. Assembly is therefore a backward patch
difference. A constant macro input already supplies its rotational
factor \(\delta\), and this patch difference supplies the second
factor in \(ob\). An affine row input gives its first-moment zero.
The other entries follow from the constant macro variation or the
constant optical patch cancellation. The actual row-supported test
has been used, rather than a bilinear product \(\chi L_\nu\).
Scalar mass inversion preserves the table: its cross blocks have one
zero, so a cross inverse followed by the \(bb\) single zero still
has the \(ob\) double zero, by \eqref{eq:mom-products}.

Laurent division by \(\tau-I=sD\) now gives
\[
 (B_W h)_b:\ Dh_b,\ h_c,\qquad
 (B_W h)_o/\delta:\ D^2h_b,\ Dh_c ,
\]
up to the stated lower coefficients. This is one derivative on the
mixed scale. A mean multiplier gives the same first-order normal
derivative estimate. An alternating \emph{velocity} multiplier
\(\delta l_c\) has two additional mesh factors relative to
\(\lambda/\delta\), and is mixed order zero.

At \(X=Y\), let \(\kappa_s=1/r+\cos\theta/\rho_\theta\).
The reference force defect is
\(R_h=KY-T_Y^*\kappa_s\). In an orthonormal Monge chart the actual
tangent-projected planar star annihilates affine nodal data exactly.
Its curved metric differs from its projected planar metric by
\(O(\delta^2)\). The projected star differs from the leading product
star by a relative \(O(\delta)\), which acts only on quadratic data.
For short/long lengths \(a_0,b_0\) and meridional length \(c_0\),
the latter star has mass row sum \(c_0(a_0+b_0)/2\); its action on
\((h_{11}x^2+2h_{12}xz+h_{22}z^2)/2\) is
\(-c_0(a_0+b_0)(h_{11}+h_{22})/2\).
Indeed its horizontal stiffness weights are \(c_0/a_0,c_0/b_0\),
its vertical weights are \((a_0+b_0)/(2c_0)\), and its diagonal
weights vanish. The affine part cancels on the actual projected
star, including its perturbation. The metric error, perturbed
quadratic weights and cubic remainder all cost \(O(\delta^3)\)
in weak force. Dividing by its \(O(\delta^2)\) mass gives
\begin{equation}\label{eq:initial-consistency}
 \norm{M^{-1}R_h}_{H^j}\leq C_j\delta,\qquad
 \norm{M^{-1}KY}_0\leq C .
\end{equation}
Each fixed coefficient difference retains the order, since the
reference descriptors are smooth periodic trigonometric functions.
Their convergent angle series give the corresponding uniform
positive analytic radius.

Since \(\Omega Y\) is tangent to the smooth torus, its weak normal
projection is \(O(\delta)\). Thus the actual initial velocity is
\[
 V(0)=-QM^{-1}R_h+\PN\varepsilon\Omega Y .
\]
It is physically \(O(\delta)\), its mean is \(O(\delta)\) in \(H^1\),
and its optical coordinate divided by \(\delta\) is uniformly bounded.
We retain this bounded initial layer.

\subsection{Force scales and the weak tangent form}

Fix an analysis constant \(\beta>0\), let \(e=X-Y\) in physical
coordinates, and define
\begin{equation}\label{eq:geometric-force}
 \widehat\kappa=(TN_{\rm ref})^{-*}N_{\rm ref}^*(KX+\beta Me),
 \qquad \widehat B=K+\beta M-H_{\widehat\kappa}.
\end{equation}
Write \(\widehat\kappa=(k_b,\lambda/\delta)\) in physical
mean/optical multiplier coordinates. A physical primal pair is
\((q_b,\delta q_c)\); pairing it with a force \((F_b,F_o)\) gives
\(F_b\cdot q_b+\delta F_o\cdot q_c\). Thus the scaled dual force is
\((F_b,\delta F_o)\), not \((F_b,F_o/\delta)\).
The inverse normal map has the state type \eqref{eq:state-type}; its force-dual type is instead
\[
 \begin{pmatrix}\text{order }0&\text{order }1\\
          \delta^2\text{ order }1&\text{order }0\end{pmatrix}.
\]
For \(F=KX+\beta Me\), the row moment calculus gives
\[
 \norm{F_b}_{H^2}+\norm{\delta F_o}_{H^3}\leq C_{\rm low},
 \qquad
 \norm{k_b}_{H^2}+\norm\lambda_{H^3}\leq C_{\rm low}.
\]
The mean force uses \(D^2y,Dw\), whereas \(\delta F_o\) uses
\(Dy,w\). The apparent extra order in the lower dual cross block
is controlled by \(\delta^2D^2\). At the reference,
\eqref{eq:initial-consistency} gives \(\lambda=O(\delta^2)\).
On an analytic geometry ball of radius \(\eta\) in \(\Xzero\),
\begin{equation}\label{eq:lambda-small}
                 \norm\lambda_{H^3}\leq C(\eta+\delta^2).
\end{equation}
All these bounds hold in the row analytic norm \eqref{eq:analytic-norm}.
The scaled normal formulas and ordered inverse series are analytic
there. Reference real-angle gaps are uniform by the same geometric
argument for every common angle shift; a small complex shift
preserves them. The holomorphic continuation of a real transpose
is used, not conjugation of the complex shift. This gives a uniform
reference series disk for the inverse coefficients.

Set \(B=\Graph^*\widehat B\Graph\), with macro/fast blocks
\(B_{ss},B_{sf},B_{fs},B_{ff}\). Their normalized classes are
\[
 \delta^{-2}\text{ class }2,\quad
 \delta^{-1}\text{ class }1,\quad
 \delta^{-1}\text{ class }1,\quad
 \text{class }0 ,
\]
and hence
\[
 B_{ss}:H^3\to H^1,\quad B_{sf}:H^2\to H^1,\quad
 B_{fs}:H^3\to H^2,\quad B_{ff}:H^2\to H^2.
\]
Products with \(\Graph\), its adjoint and global inverse kernels
preserve the factored classes by \eqref{eq:mom-products} and
\eqref{eq:adjoint-moment}. The alternating form is bounded by
\(C\norm\lambda_{H^3}\) times the mixed \(H^1\) norm squared.
For the mean part \eqref{eq:normal-derivative} gives
\[
 |\ip v{H_{k_b}v}|
     \leq C\norm{k_b}_\infty\norm v_{L^2}\norm v_{H^1}.
\]
Choose \(\beta\) large from the reference mean bound, then \(\eta\)
small relative to \(\beta\), and finally \(\delta\) small. The full
tangent form \(B\) is uniformly coercive in the mixed \(H^1\) norm.
In particular \(H:=B_{ff}\) has a uniform gap and a localized
global inverse.

The weak representation gives more information than these strong
mapping bounds. Factoring the exact reference gradient before
expanding a left adjoint gives
\begin{equation}\label{eq:gradient-factor}
       J\Graph_s=J_sD+J_0,\qquad J\Graph_f=J_f .
\end{equation}
To obtain \(J_s\), divide the optical-from-mean class \(1\) kernel
by its phase difference. Its factored zero moment belongs to \(J_0\).
The macro part uses \(D\Graph_s=\Graph_sD+[D,\Graph_s]\).
For \(J_f\), both the azimuthal gradient and \(\delta D\) are
order zero. Thus \(J_s,J_f\) have coefficient \(H^3\) from \(\Xzero\),
and \(J_0\) has \(H^2\), using \(D^2y,Dw\).

The nine stiffness coefficients are exactly
\[
 \begin{gathered}
 J_s^*J_s,\ J_s^*J_f,\ J_f^*J_s,\ J_f^*J_f,\\
 J_s^*J_0,\ J_0^*J_s,\ J_0^*J_f,\ J_f^*J_0,\ J_0^*J_0 .
 \end{gathered}
\]
For the alternating multiplier, \eqref{eq:curved-moments} places
one \(D\) on the input and one \(D^*\) on the macro test. Moving
the latter difference uses one coefficient difference. Its lower
coefficients involve \(D\lambda,D^2y,Dw,k_b\), and have \(H^2\).
The mean multiplier is a first-normal-derivative form; its
principal gradient/fast terms retain an explicit small \(\delta\).
The \(\beta\) mass is lower order, with any \(\delta D\) factor
kept inside an order-zero kernel. Consequently the exact weak blocks
can be written
\begin{equation}\label{eq:nine-blocks}
 \begin{aligned}
 B_{ss}&=D^*AD+D^*b_R+b_LD+c,\\
 B_{sf}&=D^*C_f+d_f,\qquad
 B_{fs}=C_sD+d_s,\qquad B_{ff}=H .
 \end{aligned}
\end{equation}
The principal coefficients \(A,C_f,C_s,H\) have \(H^3\);
the lower coefficients have \(H^2\) at geometry \(\Xzero\).
One additional lower-coefficient derivative costs exactly one
\(\Xone\) factor, linearly.

\begin{lemma}\label{lem:global-gap}
The actual global order-zero matrix
\(\mathbb A=\left(\begin{smallmatrix}A&C_f\\C_s&H\end{smallmatrix}\right)\)
has a uniformly positive Hermitian part on the small geometry ball.
\end{lemma}
\begin{proof}
The coherent frozen reference form is the P1 Dirichlet form on an
arbitrary macro gradient and fast coordinate. At zero phase the divided
kernel is defined by its continuous extension after removing the macro
phase zero; the class-\(1\) factorization in
\eqref{eq:gradient-factor} gives this extension. Its norm controls both:
on the product strip, vanishing gradients force the optical
coordinate and then the macro gradient to vanish. The same holds
at zero phase after division by the macro phase difference.
Compactness of phase and reference shapes gives a uniform gap.
The small optical geometry and \eqref{eq:lambda-small} preserve it;
the mean-multiplier principal terms retain their small \(\delta\).

To pass from frozen to actual kernels, take a smooth row partition
\(\sum_\nu\chi_\nu^2=1\), with bounded overlap and windows of \(L\)
rows. Kernel first moments give a commutator error \(C/L\).
Replacing coefficients in a window by its center's coherent frozen
ones costs \(CL\delta\); the inverse-kernel freezing error is
included. Therefore
\[
 \operatorname{Re}\ip v{\mathbb Av}
       \geq [c-C(L\delta+1/L)]\norm v_0^2 .
\]
Choose the integer \(L\) once, then \(\delta\) small. This proves
the global gap. The localization length is a proof choice, not
a discretization parameter.
\end{proof}

\section{A uniform analytic interval}\label{sec:interval}

The analytic weights and norms are those in \eqref{eq:analytic-norm}. The exact shifted Leibniz rule gives binomial convolution;
the polynomial factor makes this an algebra. Indeed
\((k+l+1)^p\leq C_p((k+1)^p+(l+1)^p)\), and
\(\sum(j+1)^{-2p}<\infty\). This proves the product bounds with
uniform constants. Ordered inverse series and the binomial series
for the full normal denominator converge in a sufficiently small
geometry ball. Reference angle analyticity, uniform real gaps and
complex perturbation give a positive initial radius independent
of the mesh.

\begin{proposition}\label{prop:interval}
For the fixed sufficiently small \(\varepsilon>0\), there are
\(T_0>0\), \(\rho_*>0\) and \(C\), independent of sufficiently small
\(\delta\), such that the original solution exists uniquely with
nondegenerate elements on \([0,T_0]\), and
\[
 \sup_{t\leq T_0}\left(
   \norm{\mathbf e(t)}_{\rho_*,p;\Xzero}
   +\norm{(u(t),z(t))}_{\rho_*,p;\Azero}\right)\leq C .
\]
Here \(\mathbf e=(y-y^0,w)\), the relative velocity is
\(W=V-\varepsilon\Omega X=\Graph(u,\omega)\), and
\(z=\omega+H^{-1}B_{fs}u\).
The geometry remains in a fixed small nondegenerate analytic ball.
\end{proposition}

\subsection{The relative velocity and the fast Schur variable}

The constraint is \(TW=0\). Differentiate the force equation and
use the identities
\[
 K\Omega=\Omega K,\quad M\Omega=\Omega M,\quad
 T'[\Omega X]^*\kappa=\Omega T^*\kappa .
\]
Since \((K-H_\kappa)\Omega X=-\Omega MV\), the explicit rotation
terms cancel in the full vector equation. Exactly,
\begin{equation}\label{eq:relative-equation}
 W_t=-QM^{-1}(K-H_\kappa)W+Q_tW .
\end{equation}
The normal part \(Q_tW\) is fixed by differentiating \(QW=W\).

Set \(Z=V-\beta e\). The original force equation and
\eqref{eq:geometric-force} give
\[
 \kappa=\widehat\kappa+\lambda_Z,\qquad
 \lambda_Z=R_N^*MZ=(l_b,\delta l_c),\qquad
 \norm{l_b}_{H^3}+\norm{l_c}_{H^2}
                    \leq C_{\rm low}\norm Z_{\Azero}.
\]
This is a velocity scale, distinct from the force scale
\((k_b,\lambda/\delta)\).
With \(G=\Graph^*M\Graph\) and \(\mathbf a=(u,\omega)\),
\eqref{eq:relative-equation} becomes
\[
 G\mathbf a_t+B\mathbf a
  =\beta G\mathbf a+\Graph^*H_{\lambda_Z}\Graph\mathbf a
                         -\Graph^*M\Graph'[V]\mathbf a .
\]
Equivalently, \(\mathbf a_t+G^{-1}B\mathbf a=\mathcal N_{\mathbf a}\),
where the exact source is
\begin{equation}\label{eq:physical-source}
 \mathcal N_{\mathbf a}
  =L_{\rm ref}\left[\beta W+QM^{-1}H_{\lambda_Z}W
                         +(R_N-R_M)T'[V]W\right],
 \quad R_M=M^{-1}T^*(TM^{-1}T^*)^{-1}.
\end{equation}
The moment table \eqref{eq:connection-moments}, the velocity multiplier
scale, and \eqref{eq:graph-high} give
\begin{equation}\label{eq:one-high-source}
 \norm{\mathcal N_{\mathbf a}}_{\Azero}
 \leq P(\text{low norms})
       \left(1+\norm{\mathbf e}_{\Xone}
                          +\norm{\mathbf a}_{\Aone}\right).
\end{equation}
In each product only one factor has the higher norm. The mean source
has one macro and zero optical-coordinate derivatives; the optical
source has two macro and one optical-coordinate derivatives.
These counts follow before estimation from Laurent division of
\eqref{eq:connection-moments}. Coefficient derivatives of the graph
use \(D^2y,Dw\); their one-higher applications give the single high
geometry factor in \eqref{eq:one-high-source}.

The physical mass retains its optical factor:
\[
 G=\begin{pmatrix}G_{ss}&\delta^2M_{sf}\\
             \delta^2M_{fs}&\delta^2M_{ff}\end{pmatrix}.
\]
Introduce the actual global operators
\begin{equation}\label{eq:schur-defs}
 \begin{aligned}
 C&=H^{-1}B_{fs},& z&=\omega+Cu,\\
 \overline M&=G_{ss}-\delta^2M_{sf}M_{ff}^{-1}M_{fs},&
 S&=B_{ss}-B_{sf}H^{-1}B_{fs},\\
 L&=B_{sf}-M_{sf}M_{ff}^{-1}H,&
 N&=M_{ff}^{-1}M_{fs}-C,\qquad A_f=M_{ff}^{-1}H .
 \end{aligned}
\end{equation}
Direct elimination of the mass equations, in this order, gives
\begin{equation}\label{eq:schur-system}
 \begin{aligned}
 \overline M u_t+Su+Lz&=\overline M\mathcal N_u,\\
 z_t+\delta^{-2}A_fz-N\overline M^{-1}(Su+Lz)
       &=\mathcal N_z,\\
 \mathcal N_z&=\mathcal N_\omega+C\mathcal N_u+C_tu .
 \end{aligned}
\end{equation}
The source in the original mass equation was \(G\mathcal N_{\mathbf a}\);
its mass factors cancel in this elimination. No independent source
is multiplied by \(\delta^{-2}\).
Here
\[
 C_t=H^{-1}(B_{fs}'[V]-H'[V]C),\qquad
 \widehat\kappa'[V]=R_N^*\widehat B V .
\]
Thus \(C_tu\) has the same one-high-factor estimate as
\eqref{eq:one-high-source}, without \(\kappa_t\).
The change between \((u,\omega)\) and \((u,z)\) is bounded on
\(\Azero\); its one-higher application is tame with one \(\Xone\)
factor.

\subsection{A connection for the weak macro operator}

Using \eqref{eq:nine-blocks}, define
\[
 \begin{aligned}
 a_0&=A-C_fH^{-1}C_s,& r_0&=b_R-C_fH^{-1}d_s,\\
 l_0&=b_L-d_fH^{-1}C_s,& v_0&=c-d_fH^{-1}d_s .
 \end{aligned}
\]
For arbitrary \(\xi\), test the global matrix in
Lemma~\ref{lem:global-gap} on
\((\xi,-H^{-1}C_s\xi)\). Its lower row is zero, so
\[
 \operatorname{Re}\ip{\xi}{a_0\xi}
  =\operatorname{Re}\ip{(\xi,-H^{-1}C_s\xi)}
                         {\mathbb A(\xi,-H^{-1}C_s\xi)}
       \geq c\norm\xi_0^2 .
\]
The order-zero inverse \(a_0^{-1}\) is consequently a genuine
localized global inverse. Put
\begin{equation}\label{eq:connection}
 \Gamma=a_0^{-1}r_0,\qquad \nabla=D+\Gamma .
\end{equation}
Keeping all products in their order yields the exact identity
\begin{equation}\label{eq:weak-connection}
 S=\nabla^*a_0\nabla+(l_0-\Gamma^*a_0)\nabla+(v_0-l_0\Gamma).
\end{equation}
The principal coefficient \(a_0\) has \(H^3\); \(\Gamma\) and the
two lower coefficients have \(H^2\) from geometry \(\Xzero\).
Their \(H^3\) norms cost one \(\Xone\) factor.

The choice of connection matters. For the scalar graph \((1,g)\),
the weak coefficients are \(1+g^2\), \(gg_x\), \(gg_x\), \(g_x^2\).
It gives \(\Gamma=gg_x/(1+g^2)\) and
\[
 S=\nabla^*(1+g^2)\nabla+\frac{g_x^2}{1+g^2}.
\]
The strong second derivative of the graph lies in the principal
connection, rather than in a lower potential.

For an operator coefficient, let \(A_l(k)\) map input Fourier frequency
\(k\) to \(k+l\), and define
\[
 W_p(t)^2=\sum_{j\geq0}\frac{(j+1)^{2p}t^{2j}}{(j!)^2},\qquad
 \norm A_{\rho,p;H^r,\mathrm{coeff}}^2
 =\sum_l(1+|\lambda_l|^2)^rW_p(\rho|\lambda_l|)^2
                                      \sup_k\norm{A_l(k)}^2.
\]
This is the coefficient analytic \(H^r\) norm in the next lemma.
\begin{lemma}\label{lem:same-radius}
If the actual connection has bounded coefficient analytic \(H^2\)
norm at radius \(\rho\), its common covariant jets satisfy
\[
 \sum_{j\geq0}\varpi_j^2\sum_{\ell=0}^3
                  \norm{\nabla^{j+\ell}u}_0^2
       \asymp\norm u_{\rho,p;H^3}^2 .
\]
The same assertion through offset \(2\) holds for the fast variable.
Both sides use the same radius. At exponent \(p+1/2\) the comparison
is tame: its high geometry dependence is linear, and its low
constants depend only on exponent-\(p\) norms.
\end{lemma}
\begin{proof}
The hypothesis is precisely a bound for
\(\norm\Gamma_{\rho,p;H^2,\mathrm{coeff}}\).
For a circle-dependent coefficient \(A(\vartheta)\) use
\[
 \norm A_{r,p}^2=\sum_l(1+|\lambda_l|^2)^r
                       \sup_k\norm{A_l(k,\vartheta)}_{H^p_\vartheta}^2 .
\]
Here the circle norm uses Fourier weight \((1+|m|)^{2p}\).
For \(r,p>1/2\) this is an algebra: product coefficients are
\(\sum_m A_{l-m}(k+m)B_m(k)\), and the two weighted convolution
inequalities apply in their stated order.

For \(\zeta=\rho' e^{i\vartheta}\), \(0\leq\rho'\leq\rho\), set
\[
 \Gamma_\zeta=e^{\zeta D}\Gamma e^{-\zeta D},\qquad
 \mathcal H(\zeta)=e^{\zeta\nabla}e^{-\zeta D}.
\]
The exact modular identity
\(\lambda_{k+l}-\lambda_k=e^{iks}\lambda_l\) shows that the circle
norm of the transported coefficient has precisely the factor
\(W_p(\rho'|\lambda_l|)\); the input-frequency phase just rotates
the circle. The ordered equations
\[
 \mathcal H'=\mathcal H\Gamma_\zeta,\qquad
 (\mathcal H^{-1})'=-\Gamma_\zeta\mathcal H^{-1}
\]
give bounds in the same \(r=2\) algebra by their convergent
iterated-integral series. Differentiating the interaction product
the other way gives
\[
 [D,\mathcal H]=\mathcal H\Gamma_\zeta-\Gamma\mathcal H .
\]
Its \(l,k\) coefficient contains \(e^{iks}\lambda_l\).
Thus \(\mathcal H,\mathcal H^{-1}\) gain one coefficient derivative
and belong to the \(r=3\) algebra. They act boundedly on
circle \(H^p\) with values in row \(H^3\). Circle Parseval in
\(e^{\zeta\nabla}u=\mathcal H e^{\zeta D}u\) proves same-radius
comparison for jets with an ordinary \(H^3\) base.

For each vector \(v\), the finite norm
\(\sum_{\ell=0}^3\norm{\nabla^\ell v}_0^2\) is equivalent to
\(\norm v_{H^3}^2\). Expanding \(\nabla^3-D^3\), the highest
coefficient derivative is the second derivative of \(\Gamma\).
It is row \(L^2\) and acts on an \(H^1\) field by the uniform
Fourier convolution bound. Hence
\(\norm{(\nabla^3-D^3)v}_0\leq C\norm v_{H^2}\), and triangular
induction proves the finite comparison. Apply it to
\(v=\nabla^ju\). The offset-\(2\) proof is the same finite expansion.

For \(q=p+1/2\), the algebra has the tame bound
\(\norm{AB}_{r,q}\leq C(\norm A_{r,q}\norm B_{r,p}
                     +\norm A_{r,p}\norm B_{r,q})\).
The interaction ODE and its derivative-gain identity therefore have
linear high-norm bounds, with low constants fixed at \(p\).
Since \(\Gamma\) is an analytic coefficient map using \(D^2y,Dw\),
its high norm is bounded by
\(C_{\rm low}(1+\norm{\mathbf e}_{\rho,q;\Xzero})\).
This proves the asserted tame comparison.
All series are finite-matrix entire functions before estimation,
so the coefficient integrations and Parseval passages are justified
by absolute convergence.
\end{proof}

\subsection{Common jets and the energy estimate}

The macro and fast coordinates both have dimension two. Apply the
same operator \(\nabla\) to both, using their fixed reference coordinate
identification. Put \(u_m=\nabla^m u\), \(z_m=\nabla^m z\), and
\(\mathcal T_m=[\partial_t,\nabla^m]\). Commuting the mass equation
before substituting \(u_t\) gives exactly
\begin{equation}\label{eq:jet-system}
 \begin{aligned}
 \overline M u_{m,t}+Su_m+Lz_m&=R_{u,m},\\
 z_{m,t}+\delta^{-2}A_fz_m
       -N\overline M^{-1}(Su_m+Lz_m)&=R_{z,m},
 \end{aligned}
\end{equation}
where
\begin{equation}\label{eq:jet-remainders}
 \begin{aligned}
 R_{u,m}={}&-[\nabla^m,\overline M]u_t-[\nabla^m,S]u
       -[\nabla^m,L]z+\nabla^m(\overline M\mathcal N_u)
                              +\overline M\mathcal T_m u,\\
 R_{z,m}={}&-\delta^{-2}[\nabla^m,A_f]z
       +[\nabla^m,N\overline M^{-1}S]u
       +[\nabla^m,N\overline M^{-1}L]z
                              +\nabla^m\mathcal N_z+\mathcal T_m z .
 \end{aligned}
\end{equation}
The time insertion is
\(\mathcal T_m=\sum_{k=0}^{m-1}\nabla^k\Gamma_t\nabla^{m-1-k}\).
There is no mismatch of ordinary and covariant derivatives between
the two blocks.

For a coefficient operator \(A\), the exact formulas are
\[
 [\nabla^m,A]
    =\sum_{k=1}^m\binom mk(\ad_\nabla^k A)\nabla^{m-k},
 \qquad
 e^{\zeta\nabla}Ae^{-\zeta\nabla}
       =\mathcal H(\zeta)e^{\zeta D}Ae^{-\zeta D}\mathcal H(\zeta)^{-1}.
\]
Reference positive analytic jets are \(O(\rho)\); actual-minus-reference
jets are \(O(\eta)\) on the small analytic ball. This includes the
commutator with \(\Gamma\), even for matrix coefficients. The
principal coefficients \(a_0\), masses, \(A_f\), and first-order
coefficients of \(L,N\) have analytic row \(H^3\) bounds.
The lower coefficients have \(H^2\), with one \(\Xone\) factor
for an additional derivative.
Differentiating the actual products defining \(\Gamma\) gives
\[
 \Gamma_t=a_0^{-1}(r_{0,t}-a_{0,t}\Gamma).
\]
The identity for \(\widehat\kappa'[V]\) above shows that \(\Gamma_t\)
has \(H^1\) from the velocity base \(\Azero\), and \(H^2\) with one
\(\Aone/\Xone\) factor. Low metric time derivatives are bounded by
\[
 g_V=\norm{V_b}_{W^{2,\infty}}
                       +\norm{V_o/\delta}_{W^{1,\infty}},
\]
which is controlled by \(\Azero\) and \(\Xzero\), using
\(V=\Graph(u,\omega)+\varepsilon\Omega X\).

Choose smooth real Fourier cutoffs \(\Pi_L,\Pi_H\) with
\(\Pi_L^2+\Pi_H^2=I\), low support \(|\tau-I|\leq2d_0\), and high
support \(|\tau-I|\geq d_0\). Their phase functions are fixed; their
periodized kernels have uniformly summable distance moments, by
integration by parts in the phase Fourier coefficients.
For each common jet let \(u_L=\Pi_Lu_m\), \(u_H=\Pi_Hu_m\),
and likewise for \(z\). Use the positive low energy
\[
 E_m^L=\ip{u_L}{\overline M u_L}
       +\ip{\nabla u_L}{\overline M\nabla u_L}
       +\ip{z_L}{M_{ff}z_L},
\]
and the high energy
\begin{equation}\label{eq:high-energy}
 E_m^H=\delta^{-2}\ip{\mathbf a_H}{G\mathbf a_H},
 \qquad \mathbf a_H=(u_H,z_H-Cu_H).
\end{equation}
Let \(E_m=E_m^L+E_m^H\). The current coefficients are used in these
metrics, so their equivalence constants do not depend on \(m\):
\begin{equation}\label{eq:base-energy}
 E_m\asymp\norm{u_m}_{H^1}^2+\norm{z_m}_0^2,\qquad
 D_m:=\norm{u_m}_{H^2}^2+\delta^{-2}\norm{z_m}_0^2 .
\end{equation}
For the high part this follows from physical mass positivity,
\(C:H^1\to L^2\), and
\(\delta^{-1}\norm{u_H}_0\asymp\norm{u_H}_{H^1}\).
The macro \(L^2\) addition supplies the low zero mode.

The principal low energy estimate follows from
\eqref{eq:weak-connection}. In differentiating its \(H^1\) part,
\([\nabla,\nabla^*]\) is order zero with one coefficient derivative
of \(\Gamma\). Integration by parts leaves
\(\operatorname{Re}\ip{\nabla^2u_L}{a_0\nabla^2u_L}\).
Lower terms cost a small part of \(H^2\) plus \(CE_m\).
The mass commutator contains \(u_t\); substitute the exact macro
equation in \eqref{eq:schur-system}. Its bound is
\[
 C\norm{u_L}_{H^1}
      \left(\norm{u_L}_{H^2}
                         +(d_0/\delta+1)\norm{z_L}_0\right).
\]
The \(Lz\) and \(NSu\) terms jointly have bound
\[
 C(d_0/\delta+1)\norm{u_L}_{H^2}\norm{z_L}_0,
\]
and \(NLz\) contributes
\(C(d_0^2/\delta^2+1)\norm{z_L}_0^2\).
Choose \(d_0\) small, then \(\delta\) small, and absorb these in
the two terms of \(D_m\).
On the high band, the coercivity of the physical form \(B\)
controls \(\delta^{-2}\norm{\mathbf a_H}_{\rm mixed\,H^1}^2\).
Since \(z_H=\omega_H+Cu_H\), this controls the high contribution
to \(D_m\).

We give the remainder estimates for \eqref{eq:jet-remainders},
including their high mass cross terms. The identity
\begin{equation}\label{eq:high-mass-square}
 E_m^H=\delta^{-2}\ip{u_H}{\overline M u_H}
             +\ip{z_H+Nu_H}{M_{ff}(z_H+Nu_H)}
\end{equation}
shows that macro force \(p\) and fast force \(q\) pair through
\[
 \delta^{-2}\ip{u_H}{p}
       +\ip{z_H+Nu_H}{M_{ff}(q+N\overline M^{-1}p)}.
\]
Thus the additional \(N\)-application to a macro force is retained.
For a high-cutoff physical source,
\(\norm{p_H}_0\leq C\delta\norm p_{H^1}\); its first pairing uses
the macro energy root. For force commutators, the first pairing
uses \(H^2\) dissipation. The extra \(N\)-application either has
one higher coefficient, paired with the energy root, or a
principal term bounded by a joint macro/fast dissipation.

Take \(m=j+\ell\), \(\ell=0,1,2\).
Finite reference coefficient derivatives in the stiff fast
commutator multiply lower \(\ell\) levels. The finite Leibniz
split
\(\binom{j+\ell}{k}=\sum_{r=0}^{\ell}\binom{\ell}{r}\binom{j}{k-r}\)
separates these from analytic \(j\) derivatives. Choose fixed positive
weights \(\vartheta_0=1,\vartheta_1,\vartheta_2\), successively small
relative to the lower levels, so all finite stiff Young remainders
are absorbed by their lower-level fast dissipations. Positive
analytic coefficient jets contribute \(C(\rho+\eta)\) times the
weighted dissipation. For example, the apparent factor \(j\) in a
first principal commutator is canceled by
\(\varpi_j/\varpi_{j-1}=O(\rho/j)\) against the lower dissipative jet.

For the other commutators the relevant estimates are as follows.
In \([\nabla^m,S]u\), use
\[
 [\nabla^m,S]
 =[\nabla^m,\nabla^*]a_0\nabla
   +\nabla^*[\nabla^m,a_0]\nabla
   +[\nabla^m,l_0-\Gamma^*a_0]\nabla
   +[\nabla^m,v_0-l_0\Gamma].
\]
At \(\ell=2\) the highest force coefficient is
\(D^{j+4}y,D^{j+3}w\), the same-level \(\Xzero\) geometry.
Its \(L^2\)-force pairing with \(H^2\) leaves that geometry energy
after Young's inequality. The principal positive jets have already
been absorbed. For \([\nabla^m,\overline M]u_t\), first insert
\(u_t=-\overline M^{-1}(Su+Lz)+\mathcal N_u\); its principal orders
are macro second and fast first, with the same joint low-band
estimate. The remaining mass and \(L\) coefficients have the same
geometry count.

Expand the fast \([\nabla^m,N\overline M^{-1}S]\) into its three
product commutators. Losing one principal derivative leaves at most
macro \(H^2\) at the differentiated unknown. If \(N\) differentiates
a high coefficient of \(S\), the maximum is
\(D^{j+5}y,D^{j+4}w\), one \(\Xone\) factor. It is paired with
the fast energy root. The same bound applies to
\(N\overline M^{-1}R_{u,m}\) in \eqref{eq:high-mass-square}.
The uncommuted third-order \(NSu_m\) has already been estimated
jointly in the main system.
For \([\nabla^m,N\overline M^{-1}L]z\), the remaining first fast
derivative uses the row bound \(\norm D\leq(2\delta)^{-1}\);
its principal part is absorbed by fast dissipation. Higher
coefficient terms again have one high geometry factor.

The sole explicit stiff commutator is
\(-\delta^{-2}[\nabla^m,A_f]z\).
Its analytic jets are small, and its finite reference jets are
handled by the \(\vartheta_\ell\) choice. In the high metric its
cross pairing uses
\[
 \norm{Nu_H}_0\leq C\delta\norm{u_H}_{H^2}+C\norm{u_H}_0 ,
\]
so small analytic jets also absorb the joint macro/fast terms.
The time insertions and \(\mathcal N_u,\mathcal N_z\) have one
\(\Aone\) or \(\Xone\) factor, by
\eqref{eq:physical-source}--\eqref{eq:one-high-source} and the exact
formulas for \(C_t,\Gamma_t\). They retain their energy-root pairing.

Cutoffs are applied after common jets. For every localized
coefficient \(A\), the exact kernel difference gives
\[
 \norm{[\Pi,A]}_{0\to0}
       \leq C\delta\norm{D_{\rm coeff}A}_{\rm kernel}.
\]
Apply this identity before expanding weak derivatives. Highest
macro and mass cutoff terms have an absorbable \(\delta\);
lower terms use the \(H^2\) coefficient bound with an \(H^1\)
field, or reduce \(\delta D^2\) to the available first derivative
by \(\delta\norm D\leq1/2\). Mass cutoffs still use the exact
substitution for \(u_t\). In particular the stiff high cutoff
term is exactly, up to sign and the real factor,
\[
 \delta^{-2}\ip{z_H+Nu_H}{M_{ff}[\Pi_H,A_f]z}.
\]
Its absolute value is at most
\[
 C\delta^{-1}(\norm{z_H}_0+\norm{u_H}_{H^1})\norm z_0
 \leq c_1\delta^{-2}\norm z_0^2+C_{c_1}E_m .
\]
Choose \(c_1\) as a small part of the global fast dissipation.
Base metric derivatives, including \(C_t,\Gamma_t\), use \(g_V\)
and cost only \(P(\text{low norms})E_m\).
This accounts for all mass, principal, time, source and cutoff
terms in \eqref{eq:jet-remainders}.

It remains to sum the one-higher sources without squaring a higher
norm. For a nonnegative sequence \(A=(A_j)_{j\geq0}\), let
\[
 \norm A_q^2=\sum_{j\geq0}
       \frac{(j+1)^{2q}\rho^{2j}}{(j!)^2}A_j^2 ,
\]
and define the norms of \(B,C\) in the same way. Consider
\[
 \sum_{k,l\geq0}\varpi_{k+l}^2\binom{k+l}{k}
                  A_{k+l}B_k C_{l+1}.
\]
Put \(j=k+m-1\), \(m=l+1\), and absorb weights into the sequences.
The remaining factor is exactly
\[
 \frac m\rho\,\frac{(k+m)^p}{(k+1)^p(m+1)^p}.
\]
Since \(m\leq\sqrt{j+1}\sqrt{m+1}\), weighted convolution gives
\begin{equation}\label{eq:half-pair}
 \sum_{k,l\geq0}\varpi_{k+l}^2\binom{k+l}{k}A_{k+l}B_kC_{l+1}
 \leq \frac{C_p}{\rho}
  \norm A_{p+1/2}\norm B_p\norm C_{p+1/2}.
\end{equation}
Indeed the residual polynomial kernel is bounded by
\(C_p((k+1)^{-p}+(m+1)^{-p})\), and each reciprocal sequence is
square summable. Repeated ordered products and inverse series
retain one differentiated factor. Lemma~\ref{lem:same-radius}
at half exponent converts its covariant and ordinary norms tamely.
After squaring that tame norm comparison the extra half-weight
velocity norm is bounded by a low constant times the sum of the
connection half-weight velocity energy and the ordinary half-weight
geometry energy.

Set
\[
 \begin{aligned}
 E_x&=\norm{\mathbf e}_{\rho,p;\Xzero}^2,&
 E_v&=\sum_{j\geq0}\varpi_j^2\sum_{\ell=0}^2\vartheta_\ell E_{j+\ell},\\
 D_v&=\sum_{j\geq0}\varpi_j^2\sum_{\ell=0}^2\vartheta_\ell D_{j+\ell}.
 \end{aligned}
\]
Let \(\mathcal L_x,\mathcal L_v\) be the respective energy sums
with an extra factor \(j\). In particular \(E_v\) is equivalent
to the ordinary analytic \(\Azero\) norm. In scaled mean/optical
components the geometry equation is
\[
 (y_t,w_t)=
 \left((\Graph(u,z-Cu))_b,\frac{(\Graph(u,z-Cu))_o}{\delta}\right)
                         +\varepsilon(\Omega y,\Omega w).
\]
Equation~\eqref{eq:graph-high}
has one higher factor, so \eqref{eq:half-pair} applies to it too.
The \(j=0\) part of the half weights is ordinary growth, with a
fixed constant depending on a positive lower radius.

Combining the estimates, choose \(\rho,\eta\) small to absorb
all principal \(C(\rho+\eta)D_v\) terms. For a fixed geometry
weight \(M_w\), the energy \(\mathcal E=M_wE_x+E_v\) satisfies
\begin{equation}\label{eq:combined-energy}
 \mathcal E_t+cD_v
 \leq P_{M_w}(\mathcal E)(1+\mathcal E)
  +\left(\frac{2\rho'}{\rho}
       +\frac{P_{M_w}(\mathcal E)}{\rho}\right)
                          (M_w\mathcal L_x+\mathcal L_v).
\end{equation}
The polynomial has fixed reference and positive-lower-radius
constants; none depends on the mesh.

\subsection{Continuation}

Choose the compact geometric neighborhood and \(\beta\), then the
finite \(\vartheta_\ell\) weights and \(d_0\). Choose \(\eta,\rho_0\)
small for the preceding absorptions. Initial consistency and
reference analyticity give \(E_x(0)=0\), \(E_v(0)\leq C_0\),
including the bounded scaled optical initial layer.
Choose \(M_w>(2C_0+1)/\eta^2\), and fix
\(\rho_*=\rho_0/2\).
Let \(\nu>P_{M_w}(2C_0+1)\), put \(\rho(t)=\rho_0-\nu t\),
and choose \(T_0>0\) smaller than \(\rho_0/(2\nu)\) and the
scalar comparison time for \(\mathcal E\leq2C_0+1\).
Then \eqref{eq:combined-energy} implies this energy bound and,
strictly, \(E_x<\eta^2\). This supplies the same small analytic
geometry neighborhood used in the coefficient estimates.

For every fixed mesh, all factorial jet sums and their differentiated
sums converge absolutely on compact local time intervals, since
\(D,\nabla\) and their time derivatives are finite matrices.
If the original maximal time were below \(T_0\), bounded physical
velocity would give a limit for \(X\). The scaled edge cross products
and the full weak-normal block retain strict positive gaps at that
limit. The local ODE therefore extends by Picard iteration, a
contradiction. The cyclic reduction remains the full original
solution by uniqueness. This proves Proposition~\ref{prop:interval}.

\section{An optical label and the map error}\label{sec:label}

\subsection{The label and its spatial equation}

At row \(i\), let \(p_i=Dy_i\), \(\hat e_i=p_i/|p_i|\), and define
\[
 q_i=\Omega y_i-(\Omega y_i\cdot\hat e_i)\hat e_i,\qquad
 A_i=|q_i|>0,\qquad \hat t_i=q_i/A_i,\qquad
 \hat n_i=\hat t_i\times\hat e_i.
\] These quantities are smooth
on the small neighborhood already obtained. Put
\[
 k_i=\frac{w_i\cdot\hat t_i}{A_i},\qquad
 F_i=\frac{w_i\cdot\hat n_i}{A_i}.
\]
For \(k,F\) near zero, define
\[
 d_a=(a-2k,-2F),\quad d_b=(b+2k,2F),\quad
 L_a=|d_a|,\quad L_b=|d_b|,\quad
 U=L_a+L_b,\quad V_c=L_b-L_a
\]
and the scalar label
\begin{equation}\label{eq:label}
 I(k,F)=F\left[
  \frac{256}{27(U+1)^3(1-V_c)(1+V_c)^2}
                         \right]^{1/3}.
\end{equation}
The positive cube root is smooth near zero and \(I_F(0,0)=1\).
Let \(\mathcal J(x,y)=(-y,x)\) in the \((\hat t,\hat n)\) plane.
The coherent extruded-cell normals are
\(N_a=\mathcal Jd_a/L_a\), \(N_b=\mathcal Jd_b/L_b\).
Their weighted sum is
\(\bar n=aN_a+bN_b=\mathcal J\mathcal T\), where
\(\mathcal T=ad_a/L_a+bd_b/L_b\).

Direct differentiation in \((k,F)\) verifies the label identity.
The local gradients of \(U,V_c\) are orthogonal, with
\[
 \begin{gathered}
 |\nabla_{k,F} U|^2=\frac{16(U^2-1)}{U^2-V_c^2},\qquad
 |\nabla_{k,F} V_c|^2=\frac{16(1-V_c^2)}{U^2-V_c^2},\\
 16F^2=(U^2-1)(1-V_c^2).
 \end{gathered}
\]
Moreover
\(\mathcal T=((b-a)\nabla_{k,F}U+\nabla_{k,F}V_c)/4\).
For \(F\ne0\), \eqref{eq:label} gives
\[
 \partial_U\log|I|=\frac1{U^2-1},\qquad
 \partial_{V_c}\log|I|=-\frac1{3(1-V_c^2)}.
\]
Since \(b-a=1/3\), \(\nabla_{k,F} I\cdot\mathcal T=0\).
Continuity extends this across \(F=0\). As
\(\mathcal T(0,0)=(1,0)\), there is a smooth bounded
\(\mu>0\), with bounded reciprocal, such that
\begin{equation}\label{eq:label-normal}
             \nabla_{k,F}I=\mu\bar n,\qquad
             \partial_w I_i=(\mu_i/A_i)\bar n_i .
\end{equation}
An optical component parallel to \(p_i\) has no component in this
identity, since it does not change the plane normal.

We now use the actual curved constraint, whose matrix \(T\) still
contains the original geometric normals. Write
\(W_\pm=\xi\pm\delta\zeta\). Its optical row is exactly
\[
                       T_{ob}\xi+\delta T_{oo}\zeta=0 .
\]
Whole-rectangle integration and exact divided remainders give
\[
 \sum_jT_{ob}(i,j)=\delta b_i,\qquad
 \sum_jT_{oo}(i,j)=m_i\bar n_i+\delta c_i,\qquad m_i>0 .
\]
Here \(m_i\) is a reference scalar weight bounded away from zero.
Subtract neighboring \(\xi_i,\zeta_i\) values from the finite-band
row sums, and divide by the common \(\delta\). Every difference
is \(s=4\delta\) times a finite sum of forward differences. Thus
\begin{equation}\label{eq:label-constraint}
 \bar n_i\cdot\zeta_i
  =B_i\xi_i+\sum_jC_{ij}D\xi_j
       +\delta\sum_jE_{ij}D\zeta_j+\delta G_i\zeta_i .
\end{equation}
The coefficients here are exactly the finite combinations from
\(T_{ob},T_{oo},b_i,c_i,m_i^{-1}\); they have uniform bounds and
Lipschitz dependence at geometry \(\Xzero\).
Their lower descriptors are \(D^2y,Dw\).

The position equation reads
\[
 y_t=\xi+\varepsilon\Omega y,\qquad
 w_t=\zeta+\varepsilon\Omega w,\qquad
 p_t=D\xi+\varepsilon\Omega p.
\]
The scalar \(I(y,p,w)\) is invariant under simultaneous rotation.
The explicit \(\varepsilon\) terms therefore cancel in its chain
rule. Inserting \eqref{eq:label-constraint} into
\[
 I_{i,t}=\partial_yI_i\cdot\xi_i+
        \partial_pI_i\cdot D\xi_i+
                       (\mu_i/A_i)\bar n_i\cdot\zeta_i
\]
gives an exact equation of the form
\begin{equation}\label{eq:label-spatial}
 I_t=\mathcal F_\delta(\mathbf e,\xi,D\xi)
                  +\delta\mathcal G_\delta(\mathbf e,\zeta,D\zeta).
\end{equation}
On the uniform interval, \(\mathcal F_\delta\) is uniformly locally
Lipschitz from \(\Xzero\times H^1\) to \(L^2\), and the second
term has uniformly bounded coefficient functions. This equation
retains the spatial and moving-frame terms.

\subsection{A fixed positive-time effect}

The positive lower analytic radius supplies an extra fixed row
derivative with constant \(C/\rho_*\). Proposition~\ref{prop:interval}
and \eqref{eq:graph-high} consequently give
\[
                  \norm{\mathbf e_t}_{\Xzero}\leq C,\qquad
                  \norm{\mathbf e(t)}_{\Xzero}\leq Ct .
\]
The macro equation in \eqref{eq:schur-system} gives
\(\norm{u_t}_{H^1}\leq C\) without a time derivative of \(z\).
The physical mean relative velocity is
\[
 \xi=\Graph_{bb}(\mathbf e)u
                         +\delta^2\Graph_{bo,1}(\mathbf e)\omega ,
 \qquad \omega=z-Cu ,
\]
where \(\Graph_{bo,1}\) has one row order. The second term has
size \(O(\delta^2)\) in \(H^1\), since \(\omega\) is bounded
in \(H^2\). Bounding its value at both times, rather than
differentiating it through the fast layer, gives
\[
 \norm{\xi(t)-\xi(0)}_{H^1}\leq C(t+\delta^2).
\]
Equation~\eqref{eq:label-spatial} then implies
\begin{equation}\label{eq:label-modulus}
                       \norm{I_t(t)-I_t(0)}_0\leq C(t+\delta).
\end{equation}

It remains to calculate the initial rate. At the reference,
\(I_y=I_p=0\) and \(I_w=n_0/A+O(\delta)\).
Each short or long reference quadrilateral is exactly planar.
After rotating its azimuthal midpoint to zero, its vertices are
\[
 (\rho_i\cos(d/2),\ \pm\rho_i\sin(d/2),z_i),\qquad
 (\rho_{i+1}\cos(d/2),\ \pm\rho_{i+1}\sin(d/2),z_{i+1}) .
\]
Its two geometric normals agree. The trace of
\(\varepsilon\Omega Y\) on that normal has leading magnitude
\(\varepsilon A\cos\theta\sin(d/2)\), with opposite species signs.
The short strip has the \(\chi\) sign and the long strip its negative.
For a lower optical row test on the two leading triangles, the
consistent mass gives \(A_{\rm tri}/6+A_{\rm tri}/6=A_{\rm tri}/3\);
the upper row gives the same pairing with the global \(\chi\).
After neighboring-patch assembly, the reference area factors cancel
between the optical mass and forcing. In the equation for
\(V_c=V_o/\delta\), the signed forcing-to-mass quotient is
\[
 \frac{a\sin(a\delta/2)-b\sin(b\delta/2)}{\delta(a+b)}
       =\frac{a^2-b^2}{2(a+b)}+O(\delta^2).
\]
The common triangle factor \(A_{\rm tri}/3\) cancels; the short-strip
term has positive sign and the long-strip term negative sign.
The remaining curved-row variations contribute \(O(\delta)\).
Therefore
\begin{equation}\label{eq:initial-minus-sixth}
 n_0\cdot V_c(0)
   =\varepsilon A\cos\theta\,\frac{a^2-b^2}{2(a+b)}+O(\delta)
   =-\frac{\varepsilon A\cos\theta}{6}+O(\delta).
\end{equation}
Here \(V_c=V_o/\delta\).
The term \(T_{ob}V_b(0)\), after division of the optical equation
by \(\delta\), is \(O(\delta)\), since \(T_{ob}\) is \(\delta\)
times first row order and \(V_b(0)=O(\delta)\) in \(H^1\).
Thus \(I_t(0)=-\varepsilon\cos\theta/6+O(\delta)\) in row \(L^2\).

Since \(I(0)=0\), integration of \eqref{eq:label-modulus} yields
\begin{equation}\label{eq:label-growth}
 \norm{I(t)+\varepsilon t\cos\theta/6}_0
                                    \leq C(t^2+\delta t).
\end{equation}
For sufficiently large \(n\), the exact uniform-grid norm is
\(\norm{\cos\theta}_0=\sqrt\pi\). Fix one
\[
        0<t_*<\min\{T_0,\varepsilon\sqrt\pi/(24C)\},
\]
and then take \(\delta\) small. It follows that
\begin{equation}\label{eq:positive-label}
                   \norm{I(t_*)}_0\geq
                           \varepsilon t_*\sqrt\pi/12>0 .
\end{equation}
The time is independent of the mesh and \(\varepsilon\) has already
been fixed.

\subsection{The original energy and the lifted error}

On the small neighborhood, \(|I_i|\leq C|w_i|\).
For \(E_h=X_h-Y_h\), the leading short and long azimuthal edge
differences in its aligned representation are
\[
 \delta\{a\Omega(y-y^0)-2w\},\qquad
 \delta\{b\Omega(y-y^0)+2w\}.
\]
Their two squared norms control \(w\), including every Cartesian
component. The exact rotation formulas differ by uniformly small
terms, absorbed with the physical \(L^2\) part. Reference shape
regularity therefore gives from \eqref{eq:positive-label}
\begin{equation}\label{eq:positive-H1}
             \norm{E_h(t_*)}_{H^1(\Gamma_h^0)}^2
                                      \geq c\varepsilon^2t_*^2 .
\end{equation}

Test the original force equation by \(V\) and \(\Omega X\).
The weak constraint gives \(TV=\varepsilon T\Omega X\), while
\(\ip{KX}{\Omega X}=0\) by componentwise skew rotation. Hence
the exact identity is
\begin{equation}\label{eq:original-energy}
 \frac12\frac d{dt}\ip X{KX}+\norm V_M^2
                                =\varepsilon\ip V{M\Omega X}.
\end{equation}
Here the normalized row pairings equal the full reference surface
integrals on cyclic fields, and \(\norm V_M^2=\ip V{MV}\).

Put \(d_h=\sup_{0\leq t\leq t_*}\norm{E_h(t)}_M\).
Expand \(X=Y+E_h\) in \eqref{eq:original-energy}. The work integral is
\[
 \varepsilon\ip{E_h(t)}{M\Omega Y}
          +\varepsilon\int_0^t\ip V{M\Omega E_h}\,d\tau .
\]
The reference bound \eqref{eq:initial-consistency} controls
\(\ip{E_h}{KY}\). Proposition~\ref{prop:interval} bounds
physical \(V\) uniformly in \(L^2\). Thus
\begin{equation}\label{eq:gradient-from-map}
 \frac12\norm{\nabla_{\Gamma_h^0}E_h(t)}_{L^2}^2
                   +\int_0^t\norm V_M^2\,d\tau\leq C d_h .
\end{equation}
Combining \eqref{eq:positive-H1} with this bound gives fixed
positive constants \(\alpha,C_1,C_2,D_0\), independent of the mesh,
such that
\[
 \alpha:=c\varepsilon^2t_*^2\leq C_1d_h+C_2d_h^2,\qquad
                              d_h\leq D_0 .
\]
Consequently \(d_h\geq\alpha/(C_1+C_2D_0)>0\).
Lemma~\ref{lem:admissible} gives a fixed lift norm constant \(c_\ell>0\)
and an \(O(\delta^2)\) identity error, so
\[
 \sup_{t\leq t_*}\norm{X_h^\ell(t)-\id}_{L^2(\Gamma^0)}
           \geq \frac{c_\ell\alpha}{C_1+C_2D_0}-C\delta^2 .
\]
For sufficiently small \(\delta\), take
\(T_*=t_*\) and
\(c_*=c_\ell\alpha/[2(C_1+C_2D_0)]\).
These are positive fixed constants. The continuous solution is
the identity, the closest lift is the original one, and the full
original discrete solution exists nondegenerately on this interval.
This proves Theorem~\ref{thm:failure}.

\section{Convergence for cubic elements}\label{sec:cubic}

We now consider arbitrary smooth data, rather than the stationary torus.
Write \(u=U_{\rm ext}\) in this section. Assume that
\eqref{eq:continuous} has a smooth solution on a fixed finite interval
\([0,T]\), with smooth flow map and inverse throughout that interval.
Let \(\Gamma_h^0\) be any admissible shape-regular, quasi-uniform cubic
nodal interpolant of \(\Gamma^0\). For its element parametrizations require
exactly the original bounds
\begin{equation}\label{eq:cubic-mesh}
 \max_K\bigl(\norm{F_K}_{W^{3,\infty}(K_f)}
       +\norm{F_K^{-1}}_{W^{1,\infty}(K)}\bigr)\leq C_0 .
\end{equation}
Here \(K_f\) is the flat triangle with the same vertices as \(K\).
The space \(S_h\) now consists of continuous functions whose pullbacks
by \(F_K\) are cubic polynomials. We retain precisely
\eqref{eq:discrete-time}--\eqref{eq:discrete-force}, including their
fixed reference measure and discontinuous current unit normals.

\begin{theorem}\label{thm:cubic}
For every such smooth solution and mesh family there are \(h_0,C_T>0\),
independent of \(h\), such that for \(h<h_0\) the original semidiscrete
solution exists uniquely on \([0,T]\). Its map remains an embedding and
its element differentials remain nondegenerate. Moreover,
\begin{equation}\label{eq:cubic-result}
 \sup_{0\leq t\leq T}
 \norm{X_h^\ell(t)-X(t)}_{L^2(\Gamma^0)}
       \leq C_T h^3\sqrt{1+|\log h|}.
\end{equation}
If \(d(x,t)\) is signed distance to \(\Gamma(t)\), then also
\[
 \sup_{0\leq t\leq T}\norm{d(X_h(\cdot,t),t)}_{L^2(\Gamma_h^0)}
       \leq C_T h^3 .
\]
The constants may depend on the fixed smooth solution, \(T\), and the
uniform mesh-family bounds.
\end{theorem}

The proof controls shape error before estimating the full map error.
The first energy is a relative volume weighted by signed distance.
The second is the energy of a weakly tangent part of the map error.
Their interaction leaves a vanishing coefficient in the shape equation,
while allowing a logarithmic loss in the tangent estimate.

\subsection{Reference geometry and weak tangency}

All notation in this section is local; in particular \(D\) below is a
gradient norm, rather than the row difference of the preceding sections.
Put
\[
 (w,z)_h=\int_{\Gamma_h^0}w\cdot z,\qquad
 a_h(w,z)=\int_{\Gamma_h^0}\nabla w:\nabla z,
 \qquad (L_hw,z)_h=a_h(w,z).
\]
Unqualified norms and gradients use this fixed surface. Let \(\Pi_h\)
and \(\Pi_h^3\) be the scalar and vector mass projections. With the
initial closest-point projection \(a:\Gamma_h^0\to\Gamma^0\), define
\[
 \overline X=X\circ a,\quad Y=I_h\overline X,\quad Z_h=Y_t,
 \quad m=n(\overline X,t),\quad\nu=N(Y),\quad
 \kappa^*=I_h(\kappa\circ a).
\]
The known unit field \(m\) is the restriction of a smooth field on a
fixed initial tube. It is used only in the analysis.

The original mesh assumptions imply
\begin{equation}\label{eq:cubic-interpolation}
 \norm{Y-\overline X}_\infty\leq Ch^4,\qquad
 \norm{\nabla(Y-\overline X)}_\infty
       +\norm{\nu-m}_\infty\leq Ch^3.
\end{equation}
Indeed, derivatives of \(F_K\) through order three are bounded and
its fourth derivatives vanish. Composing any fixed smooth ambient
function with \(F_K\) therefore gives bounded fourth derivatives.
Subtract its cubic Taylor polynomial and use the fixed nodal basis
on the scaled triangle. This gives interpolation errors \(Ch^4\)
and \(Ch^3\) in value and gradient. Applied to the initial signed
distance, which vanishes at all nodes, it gives initial distance
\(Ch^4\) and initial normal error \(Ch^3\).
The closest-point area ratio is \(1+O(h^4)\), and the pulled-back
Dirichlet metric differs by \(O(h^4)\): its displacement contribution
is \(O(h^4)\), and its squared normal-slope contribution is \(O(h^6)\).

Adjacent elements have the same polynomial edge tangent. Their outward
conormals therefore have sum \(O(h^3)\). For a known smooth ambient
coefficient \(b\), this gives
\begin{equation}\label{eq:cubic-conormal}
                    [\partial_\mu b]=O(h^3).
\end{equation}
Scaled polynomial norm equivalence also gives
\begin{equation}\label{eq:cubic-inverse}
 \begin{gathered}
 \norm w_\infty\leq Ch^{-1}\norm w_2,\qquad
 \norm{\nabla w}_2\leq Ch^{-1}\norm w_2,\\
 \norm{\nabla w}_\infty\leq Ch^{-2}\norm w_2,\qquad
 \norm w_{L^2(\text{edges})}\leq Ch^{-1/2}\norm w_2.
 \end{gathered}
\end{equation}
The last estimate is summed over element traces, with bounded overlap.
In particular \(\norm{L_hw}_2\leq Ch^{-1}\norm{\nabla w}_2\).
These geometric estimates concern the known initial mesh, not the
curvature of the unknown evolving surface.

For a smooth coefficient \(b\) and an FE function \(w\), product
superapproximation takes the form
\begin{equation}\label{eq:cubic-product}
 \norm{bw-I_h(bw)}_2+h\norm{\nabla(bw-I_h(bw))}_2
                          \leq Ch^2\norm w_{H^1}.
\end{equation}
To see the two powers, remove a constant value of \(b\), whose product
with \(w\) is reproduced. Subtract also the element mean of \(w\).
The latter difference has norm at most \(Ch\norm{\nabla w}_2\), by
the line-segment formula on the flat triangle. In the fourth-order
interpolation remainder, the highest remaining polynomial derivative
is bounded by \(Ch^{-2}\norm{\nabla w}_2\), so multiplication by
\(h^4\) gives \(h^2\). The constant part of \(w\) gives
\(Ch^4\norm w_2\). Differentiating the scaled basis once proves the
gradient assertion. Thus
\begin{equation}\label{eq:cubic-project-product}
 \norm{(I-\Pi_h^3)(m\chi)}_2\leq Ch\norm\chi_2,
 \qquad\norm{\Pi_h^3(m\chi)-m\chi}_{H^1}\leq C\norm\chi_2.
\end{equation}
For the second bound split at \(I_h(m\chi)\) and use
\eqref{eq:cubic-inverse}; no \(H^1\) stability of \(\Pi_h\) is needed.

For any measurable unit field \(b\), set
\[
 \mathcal T_b v=\Pi_h(b\cdot v),\qquad
 \mathcal T_b^*\chi=\Pi_h^3(b\chi),\qquad
 \mathcal S_b=\mathcal T_b\mathcal T_b^*.
\]
Mass orthogonality gives the exact identity
\begin{equation}\label{eq:cubic-schur-identity}
 (\mathcal S_b\chi,\chi)_h
       =\norm\chi_2^2-\norm{(I-\Pi_h^3)(b\chi)}_2^2.
\end{equation}
Hence \(\mathcal S_m\geq1-Ch^2\). Write
\[
 \mathcal T=\mathcal T_m,\quad\mathcal S=\mathcal S_m,\quad
 Q=I-\mathcal T^*\mathcal S^{-1}\mathcal T,
 \quad e=X_h-Y=z+\mathcal T^*\psi,
\]
\[
 z=Qe,\quad\psi=\mathcal S^{-1}\mathcal T e,\qquad
 E=\norm e_2,\quad F=\norm\psi_2,\quad Z=\norm z_2,
 \quad D=\norm{\nabla z}_2.
\]
Then \(Q\) is the mass-orthogonal projection onto \(\ker\mathcal T\),
\begin{equation}\label{eq:cubic-split}
 E^2=Z^2+(\mathcal S\psi,\psi)_h\asymp Z^2+F^2,
 \qquad (e_t,z)_h=\tfrac12(Z^2)'-\tfrac12(e,Q_t e)_h,
 \qquad\norm{Q_t}_{2\to2}\leq C.
\end{equation}
The time bound follows by differentiating the inverse and using
\(\mathcal T_t v=\Pi_h(m_t\cdot v)\).

Weak tangency yields a stronger leakage estimate than an inverse bound:
\begin{equation}\label{eq:cubic-leak}
 \norm{m\cdot z}_2\leq Ch^2(D+Z),\qquad
 \norm{\nabla(m\cdot z)}_2\leq Ch(D+Z).
\end{equation}
Indeed, \(\Pi_h(m\cdot z)=0\), so
\(I_h(m\cdot z)=\Pi_h[I_h(m\cdot z)-m\cdot z]\).
Apply \eqref{eq:cubic-product} and then an inverse estimate.
Furthermore,
\begin{equation}\label{eq:cubic-cross}
 \begin{gathered}
 |a_h(z,\mathcal T^*\chi)|\leq C(D+Z)\norm\chi_2,
 \quad\norm{\mathcal T L_hz}_2\leq C(D+Z),\\
 \norm{QL_h\mathcal T^*\psi}_2\leq Ch^{-1}F,
 \qquad a_h(e,z)\geq\tfrac12D^2-CE^2.
 \end{gathered}
\end{equation}
For the first assertion, replace \(\mathcal T^*\chi\) by \(m\chi\)
using \eqref{eq:cubic-project-product}, and expand
\[
 a_h(z,m\chi)=(\nabla\chi,\nabla(m\cdot z))_h
 -\sum_j(\nabla\chi,z_j\nabla m_j)_h
 +\sum_j(\chi\nabla m_j,\nabla z_j)_h.
\]
The first term uses \eqref{eq:cubic-leak}. Integrate the middle term
on each initial element; the volume coefficients are bounded, and
\eqref{eq:cubic-conormal} with two inverse traces bounds its edge term
by \(Ch^2\norm\chi_2Z\). The last term is bounded directly.
Adjoint testing gives the third assertion of \eqref{eq:cubic-cross};
Young's inequality gives the last.

We also use the two-dimensional discrete logarithmic inequality
\begin{equation}\label{eq:cubic-log}
 \norm z_\infty\leq C\sqrt\ell(D+Z),\qquad\ell=1+|\log h|.
\end{equation}
Here is a proof uniform over these curved meshes. The element mass
matrix gives a polynomial evaluation dual \(\delta_h\), with
\(\norm{\delta_h}_1\leq C\), \(\norm{\delta_h}_2\leq Ch^{-1}\).
Push its density into a smooth initial coordinate chart. With
\(\widehat f(\xi)=\int f(x)e^{-2\pi i x\cdot\xi}\,dx\), Plancherel's
identity~\cite[Theorem~43]{TaoFourier} and \(|\widehat\delta_h|\leq C\)
give
\[
 \int_{\R^2}\frac{|\widehat\delta_h(\xi)|^2}{1+4\pi^2|\xi|^2}\,d\xi
 \leq C\int_{|\xi|\leq h^{-1}}\frac{d\xi}{1+|\xi|^2}
             +Ch^2\norm{\delta_h}_2^2\leq C\ell.
\]
Weighted Cauchy--Schwarz bounds its pairing with a chart \(H^1\)
function. The Fourier derivative identity follows by integration by
parts and then convolution approximation. A fixed finite chart cover,
with cutoffs equal to one on the relevant elements, proves
\eqref{eq:cubic-log}.

\subsection{The actual mixed system on a small-error interval}

Fix the analysis exponent \(1/4\), and work initially while
\begin{equation}\label{eq:cubic-first-bootstrap}
 E\leq h^{11/4},\qquad
 \rho=C_1(h^3+h^{-2}E),\qquad\alpha=h+\rho .
\end{equation}
The fixed \(C_1\) is sufficiently large. For small \(h\),
\(\rho\leq Ch^{3/4}\), \(\alpha\geq h\),
\(\alpha^2\leq h\), and \(\alpha^2\leq Ch^{3/2}\).
All constants below are uniform on one fixed small geometric
neighborhood and independent of the later bulk constant.

For the actual normal \(N=N(X_h)\),
\begin{equation}\label{eq:cubic-normal-difference}
 \norm{N-m}_\infty\leq C\rho,\qquad
 \norm{N-m}_2\leq C(h^{-1}F+E+\alpha D+h^3).
\end{equation}
The first variation about \(Y\) is \(-\mathcal H(\mathrm De)^{\mathsf T}\nu\),
where \(\mathcal H\) is a uniformly bounded inverse-tangent coefficient.
Use
\((\mathrm De)^{\mathsf T}m=\nabla(m\cdot e)-(\mathrm Dm)^{\mathsf T}e\),
\eqref{eq:cubic-leak}, and
\(\norm{\nabla e}_2\leq C(D+h^{-1}F)\).
The linear term is bounded by \(C(h^{-1}F+E+hD)\); its Taylor remainder
is at most \(C\rho(D+h^{-1}F)\). Adding the reference normal error
proves \eqref{eq:cubic-normal-difference}.

Set \(\mathcal U=\mathcal T_N\) and
\(Q_N=I-\mathcal U^*\mathcal S_N^{-1}\mathcal U\).
The exact identity \eqref{eq:cubic-schur-identity} gives
\begin{equation}\label{eq:cubic-actual-schur}
 \mathcal S_N\geq[1-(Ch+C\rho)^2]I,
 \qquad\norm{Q_N-Q}_{2\to2}\leq C\rho .
\end{equation}
The inverse-difference formula proves the second bound. Discontinuity
of \(N\) across edges is permitted; no derivative of \(N\) is used.

Define the known defects
\[
 \mathfrak d=Z_h+L_hY-\mathcal T_\nu^*\kappa^*,\qquad
 r_h=\Pi_h[(Z_h-u(Y,t))\cdot\nu].
\]
Pulling the smooth weak equation to the initial mesh and using
\eqref{eq:cubic-interpolation} gives
\begin{equation}\label{eq:cubic-defects}
 \norm{\mathfrak d}_{H_h^{-1}}\leq Ch^3,
 \qquad\norm{\mathfrak d}_2\leq Ch^2,
 \qquad\norm{r_h}_2\leq Ch^3.
\end{equation}
The dual norm here tests only FE vectors in their full \(H^1\) norm.
The stiffness interpolation error is \(Ch^3\), while initial metric
errors are \(O(h^4)\). Inverse testing gives the mass bound. For the
bottom defect compare with \(\overline X_t-u(\overline X)\), which is
smooth and exactly \(m\)-tangent.

For \(\mu=\kappa_h-\kappa^*\), subtraction gives exactly
\begin{equation}\label{eq:cubic-error}
 e_t+L_he=\mathcal U^*\mu+G,\qquad\mathcal Ue_t=R,
\end{equation}
where
\[
 \begin{split}
 G&=\Pi_h^3[\kappa^*(N-\nu)]-\mathfrak d,\\
 R&=\Pi_h[(u(X_h)-u(Y))\cdot N
               +(u(Y)-Z_h)\cdot(N-\nu)]-r_h .
 \end{split}
\]
Thus \(\norm G_2+\norm R_2\leq C(h^{-1}F+E+\alpha D+h^2)\),
with \(h^3\) in the separate bound for \(R\).
Exact elimination gives
\[
 \mu=\mathcal S_N^{-1}(R+\mathcal UL_he-\mathcal UG),\qquad
 e_t=-Q_NL_he+Q_NG+\mathcal U^*\mathcal S_N^{-1}R .
\]
By \eqref{eq:cubic-cross} and
\(\norm{\mathcal U-\mathcal T}\leq C\rho\),
\begin{equation}\label{eq:cubic-mixed-bounds}
 \begin{split}
 \norm\mu_2&\leq C(h^{-2}F+\alpha h^{-1}D+E+h^2),\\
 \norm{e_t}_2&\leq C(h^{-1}D+\alpha h^{-2}F+E+h^2).
 \end{split}
\end{equation}
For the improved normal coefficient in the second line, keep the
projection before the stiffness:
\[
 \norm{Q_NL_h\mathcal T^*\psi}_2
       \leq C(h^{-1}+\rho h^{-2})F.
\]

Finally, the actual pointwise residual
\(q=(v_h-u(X_h,t))\cdot N\) satisfies \(\Pi_hq=0\) and
\begin{equation}\label{eq:cubic-q}
 \norm q_2\leq C(h^{-1}F+\alpha h^{-1}D+E+\alpha h^2).
\end{equation}
Indeed, its \(e_t\cdot m\) projection defect costs \(Ch\norm{e_t}_2\),
and its \(e_t\cdot(N-m)\) part costs \(C\rho\norm{e_t}_2\).
The remaining reference-tangent part costs
\(C(\norm{N-m}_2+E+h^4)\). Substitution of
\eqref{eq:cubic-mixed-bounds} gives an additional
\(\alpha^2h^{-2}F\), absorbed by \(\alpha^2\leq h\);
\(h^3\leq\alpha h^2\) handles the source.

\subsection{Relative volume and the area-weighted projection defect}

Orient one connected component outward, and let \(\Omega(t)\) and
\(\Omega_h(t)\) be its exact and numerical bounded interiors. The
small geometric neighborhood gives a normal graph \(r(s,t)\) over
\(\Gamma(t)\). Define
\begin{equation}\label{eq:cubic-bulk}
 \mathcal B(t)=\int_{\R^3}
       (\mathbf1_{\Omega_h(t)}-\mathbf1_{\Omega(t)})d(x,t)\,dx .
\end{equation}
Only a fixed tube is involved, so signed distance may be smoothly
extended outside it. In normal coordinates,
\[
 \mathcal B=\int_{\Gamma(t)}\int_0^{r(s,t)}
       \xi\det(I+\xi\mathsf S(s,t))\,d\xi\,dS_s,
\]
where \(\mathsf S\) is the exact shape operator, with the sign matching
this normal-coordinate Jacobian. The integral is nonnegative also
when \(r<0\). Uniform graph and area bounds give, for \(f=d(X_h,t)\),
\begin{equation}\label{eq:cubic-bulk-distance}
              c\norm f_2^2\leq\mathcal B\leq C\norm f_2^2.
\end{equation}
For several components take the sum of these functionals separately,
so nesting causes no cancellation. Reversing a component's orientation
reverses its multiplier and leaves the position equations unchanged.

Let \(J\) be the area ratio of \(X_h\) from the fixed initial mesh to
the current numerical surface. Closed-surface transport gives
\begin{equation}\label{eq:cubic-bulk-transport}
 \mathcal B'=\int(\mathbf1_{\Omega_h}-\mathbf1_\Omega)
             [d_t+\operatorname{div}(du)]\,dx+(q,fJ)_h.
\end{equation}
To verify the formula for piecewise-smooth elements, differentiate
their swept regions. Each face contributes its area times normal
displacement; the common moving-edge traces cancel all side fluxes.
Add and subtract \(u\) on the two closed boundaries and apply the
divergence identity to \(du\). The exact normal speed equals
\(u\cdot n\), giving \eqref{eq:cubic-bulk-transport}. No derivative
of the numerical normal appears.

On the exact surface, \(d_t+u\cdot\nabla d=0\). Integrating its smooth
normal derivative along tube fibers shows that this expression is
\(c_1(x,t)d\) with bounded \(c_1\). The volume term in
\eqref{eq:cubic-bulk-transport} is therefore at most \(C\mathcal B\).
The ORIGINAL weak constraint, in the fixed measure, removes only an
FE projection:
\begin{equation}\label{eq:cubic-bulk-weak}
 \mathcal B'\leq C\mathcal B+(q,(I-\Pi_h)(fJ))_h.
\end{equation}
Thus the current area ratio \(J\) must remain in the test defect.

The estimate that permits closure is
\begin{equation}\label{eq:cubic-projection-defect}
 \norm{(I-\Pi_h)(fJ)}_2
 \leq C\{(h+\rho^2)F+h\alpha(D+E)+h^4\}
 \leq C\{hF+h\alpha(D+E)+h^4\}.
\end{equation}
We give all parts of its proof. Let \(j\) be the known smooth exact
flow area ratio, pulled back to the initial mesh. Distance Taylor
expansion at \(\overline X\) gives
\begin{equation}\label{eq:cubic-distance-expansion}
 f=m\cdot e+\mathcal R_f,\qquad
 \norm{\mathcal R_f}_2\leq C(h^4+h^{-1}E^2),\qquad
 \norm f_\infty\leq C(h^4+h^{-1}E).
\end{equation}
Use the weak leakage of \(z\),
\(\norm{m\cdot\Pi_h^3(m\psi)-\psi}_2\leq ChF\), and the
smooth-product defect of \(j\psi\). They imply
\begin{equation}\label{eq:cubic-smooth-area}
 \norm{(I-\Pi_h)(jf)}_2
      \leq C\{hF+h^2(D+E)+h^4+h^{-1}E^2\}.
\end{equation}

For a rank-two differential matrix \(A\), write
\(\mathcal J(A)=\sqrt{\det(A^{\mathsf T}A)}\). Its derivative at the
exact differential satisfies
\begin{equation}\label{eq:cubic-area-cancellation}
 \mathrm D\mathcal J(\mathrm D\overline X)[m\otimes\xi]=0,
\end{equation}
since \((\mathrm D\overline X)^{\mathsf T}m=0\).
Write \(e=z+m\psi+r\), where
\(r=\Pi_h^3(m\psi)-m\psi\) has \(H^1\) norm at most \(CF\).
The term \(m\nabla\psi\) therefore drops out of the linear area
variation. Its other linear terms cost \(C(D+E)\), and its quadratic
remainder costs \(C\rho(D+h^{-1}F)\). Interpolation and initial lift
errors give
\begin{equation}\label{eq:cubic-area-difference}
 \norm{J-j}_2\leq C\{h^3+D+E+\rho(D+h^{-1}F)\}.
\end{equation}
Multiply by the last bound in \eqref{eq:cubic-distance-expansion}.
The inequality \(E/h\leq Ch\rho\) gives
\[
 \norm{f(J-j)}_2
       \leq C\{(h+\rho^2)F+h\alpha(D+E)+h^4\}.
\]
Keep the term \(h^{-1}E^2\) in \eqref{eq:cubic-smooth-area} until using
\(h^{-1}E^2\leq Ch\rho E\). Combining proves the first inequality
in \eqref{eq:cubic-projection-defect}; \(\rho^2\leq\alpha^2\leq h\)
proves the second.

The same decomposition, absorbing \(ChF\), gives
\[
 F\leq C\norm f_2+Ch^2(D+E)+Ch^{-1}E^2+Ch^4.
\]
Since \(D\leq Ch^{-1}E\),
\begin{equation}\label{eq:cubic-shape-normal}
 F^2\leq C\mathcal B+Ch^2E^2+Ch^{-2}E^4+Ch^8.
\end{equation}
Pair \eqref{eq:cubic-q} and \eqref{eq:cubic-projection-defect} in
\eqref{eq:cubic-bulk-weak}. Their main terms are \(F^2\),
\(\alpha FD\), and \(\alpha^2D^2\). All cross terms are bounded by
\(C\{F^2+\alpha^2(D^2+E^2)+h^6\}\). In particular the source
\(h^4E\) costs
\(\alpha^2E^2+h^8/\alpha^2\leq\alpha^2E^2+h^6\);
the other source pairings use Young's inequality in the same way.
Under \eqref{eq:cubic-first-bootstrap},
\(h^{-2}E^4\leq h^{7/2}E^2\leq h^2E^2\). Thus
\begin{equation}\label{eq:cubic-bulk-energy}
 \mathcal B'\leq C\mathcal B+C\alpha^2(D^2+E^2)+Ch^6.
\end{equation}
The vanishing coefficient \(\alpha^2\) permits control of shape error
without first knowing an optimal bound for the whole map.

\subsection{Tangent energy and the fourth power of normal error}

Test \eqref{eq:cubic-error} by \(z\). Equations
\eqref{eq:cubic-split} and \eqref{eq:cubic-cross} give
\begin{equation}\label{eq:cubic-tangent-start}
 \tfrac12(Z^2)'+\tfrac12D^2
       \leq CE^2+(G,z)_h+(\mu,(\mathcal U-\mathcal T)z)_h.
\end{equation}
For any fixed \(\vartheta>0\), small \(h\) gives
\begin{equation}\label{eq:cubic-G}
             |(G,z)_h|\leq\vartheta D^2+C_\vartheta(E^2+h^6).
\end{equation}
The defect \(\mathfrak d\) uses its \(H_h^{-1}\) bound. For the linear
normal variation write
\[
 -\mathcal H\nabla(e\cdot m)+\mathcal H(\mathrm Dm)^{\mathsf T}e
       -\mathcal H(\mathrm De)^{\mathsf T}(\nu-m).
\]
Replacing \(\mathcal H\) by the exact-reference coefficient costs
\(Ch^3\norm{\nabla e}_2E\leq Ch^2E^2\). The replacement of
\(\kappa^*\) by its smooth reference value is smaller. Integrate the
first term once on each initial element. Volume terms cost
\(CE(D+E)\). The known coefficient flux jumps are \(O(h^3)\), since
their dependence on the initial tangent planes is smooth and their
conormals have the jump in \eqref{eq:cubic-conormal}. Two inverse
traces bound the edge terms by \(CE^2\).

For the nonlinear remainder, choose an orthonormal reference tangent
frame \(\mathsf V\) at \(Y\), and write \(\mathrm DY=\mathsf V R_0\).
Both \(R_0\) and its inverse are uniformly bounded. Normalize columns by
\(R_0^{-1}\), so
\[
 \mathrm DX_hR_0^{-1}=\mathsf V A+\nu b^{\mathsf T},\qquad
 N=\frac{\nu-\mathsf V A^{-\mathsf T}b}
                {\sqrt{1+|A^{-\mathsf T}b|^2}}.
\]
Here \(A-I\) is the tangential differential perturbation and \(b\)
the normal slope. After subtracting the first variation, the remainder
\(\mathcal R\) satisfies
\begin{equation}\label{eq:cubic-refined-remainder}
 |P_\nu\mathcal R|\leq C(|A-I||b|+|b|^3),\qquad
 |\nu\cdot\mathcal R|\leq C|b|^2,
 \quad P_\nu=I-\nu\nu^{\mathsf T}.
\end{equation}
Tangential projection removes \(m\nabla\psi\) up to an
\(O(h^3)\nabla\psi\) error. Hence
\[
 \norm{A-I}_2\leq C(D+E),\quad\norm b_\infty\leq C\rho,
 \quad\norm b_2\leq C(D+h^{-1}F),\quad
 \norm{\nu\cdot z}_2\leq Ch^2(D+E)+Ch^3E.
\]
The tangential remainder pairing is bounded by
\(C\rho E(D+E)+C\rho^2E(D+h^{-1}F)\); the normal pairing is at most
\(C\rho(D+h^{-1}F)[h^2(D+E)+h^3E]\).
Use \(F\leq CE\) and \(\rho^2/h\leq Ch^{1/2}\to0\).
Young's inequality gives \eqref{eq:cubic-G}.

The multiplier term retains a particular normal-error factor. From
\eqref{eq:cubic-normal-difference}, \eqref{eq:cubic-mixed-bounds},
and \eqref{eq:cubic-log},
\begin{equation}\label{eq:cubic-multiplier-product}
 \begin{split}
 |(\mu,(\mathcal U-\mathcal T)z)_h|
 &\leq C\sqrt\ell(h^{-2}F+\alpha h^{-1}D+E+h^2)\\
 &\hspace{12mm}\cdot(h^{-1}F+E+\alpha D+h^3)(D+E).
 \end{split}
\end{equation}
Its leading product is bounded by
\begin{equation}\label{eq:cubic-critical}
 C\sqrt\ell h^{-3}F^2(D+E)
    \leq\vartheta D^2+CE^2+C_\vartheta\ell h^{-6}F^4.
\end{equation}
To account for every other term, the product of the first two
parentheses in \eqref{eq:cubic-multiplier-product}, less
\(h^{-3}F^2\), is exactly
\[
 \begin{split}
 &(h^{-2}+h^{-1})FE+2\alpha h^{-2}FD
 +(\alpha h^{-1}+\alpha)ED+\alpha^2h^{-1}D^2+E^2\\
 &\hspace{15mm}+2hF+2\alpha h^2D+(h^2+h^3)E+h^5.
 \end{split}
\]
After multiplication by \(D+E\), the first eight groups have respective
bounds by \(D^2+E^2\) with coefficients
\[
 C h^{3/4},\quad C h^{3/2},\quad C h^{5/2},\quad C h^{9/4},
 \quad C h^{11/4},\quad Ch,\quad C h^{11/4},\quad Ch^2.
\]
For the fourth group the bound uses only \(D^2\). These follow from
\(F\leq CE\), \(D\leq Ch^{-1}E\), and
\eqref{eq:cubic-first-bootstrap}. Each coefficient still tends to zero
after multiplication by \(\sqrt\ell\). The final source gives
\(\vartheta D^2+CE^2+C_\vartheta h^{10}\ell\). Absorbing the
small terms in \eqref{eq:cubic-tangent-start} proves
\begin{equation}\label{eq:cubic-tangent-energy}
                 (Z^2)'+cD^2\leq CE^2+Ch^6+C\ell h^{-6}F^4,
                 \qquad c>0.
\end{equation}
The fourth power here is \(F^4\), rather than \(E^4\).

\subsection{Closing the estimates and continuing to the final time}

Add a second bootstrap \(\mathcal B\leq Mh^6\), with \(M\) fixed
independently of \(h\). At time zero \(e=0\), since all initial nodes
are on \(\Gamma^0\), and \(\mathcal B(0)\leq Ch^8\). Thus both
bootstraps hold strictly on a nonempty initial interval. By
\eqref{eq:cubic-shape-normal},
\[
 F^2\leq CMh^6+Ch^{15/2}+Ch^9+Ch^8\leq C(M+1)h^6.
\]
Insert this into \eqref{eq:cubic-tangent-energy} and use
\eqref{eq:cubic-split}. The integrating factor \(e^{-Ct}\) yields
\begin{equation}\label{eq:cubic-map-energy}
 \sup_{s\leq t}E(s)^2+\int_0^tD(s)^2\,ds
                              \leq C_{M,T}h^6\ell.
\end{equation}
Then \eqref{eq:cubic-bulk-energy} and \(\alpha^2\leq Ch^{3/2}\) give
\begin{equation}\label{eq:cubic-bulk-improvement}
       \mathcal B(t)\leq C_T h^6+C_{M,T}h^{3/2}h^6\ell.
\end{equation}
The first \(C_T\) is independent of \(M\): it comes from the linear
bulk equation, the initial bulk, and the smooth \(Ch^6\) source.
Choose \(M\geq4C_T\) first, and then choose \(h\) small enough that
the second term is at most \(C_Th^6\). This gives
\(\mathcal B\leq Mh^6/2\). Also
\eqref{eq:cubic-map-energy} gives \(E\leq h^{11/4}/2\), since
\(h^{1/4}\sqrt\ell\to0\). Both bootstraps improve strictly.

For completeness, full-time existence is part of this argument.
Where element cross products are nonzero and \(\mathcal S_N\) is
invertible, the ORIGINAL system is the smooth finite-dimensional ODE
\begin{equation}\label{eq:cubic-ODE}
 \begin{split}
 X_{h,t}&=-Q_NL_hX_h+
         \mathcal U^*\mathcal S_N^{-1}\Pi_h[N\cdot u(X_h,t)],\\
 \kappa_h&=\mathcal S_N^{-1}
             \{\Pi_h[N\cdot u(X_h,t)]+\mathcal UL_hX_h\}.
 \end{split}
\end{equation}
The initial Schur bound gives local existence and uniqueness: on a small
fixed-\(h\) nodal ball the vector field has finite value and derivative
bounds, and successive integral iterates have factorially summable
differences. The same difference estimate proves uniqueness.
No Lipschitz constant uniform in \(h\) is required.

On its maximal safe interval, the improved estimates imply
\[
 \norm{X_h-\overline X}_{W^{1,\infty}}
       \leq C(h^3+h^{-2}E)\leq C_T h\sqrt\ell\longrightarrow0.
\]
The first bootstrap alone already bounds this by \(Ch^{3/4}\).
The exact embedding has a uniform differential lower bound and local
graph radius on compact space-time. On each smaller graph chart,
project to its tangent plane and integrate the derivative difference
along line segments. Small piecewise \(W^{1,\infty}\) perturbations
preserve a local lower Lipschitz bound. Continuous FE traces make
these restrictions absolutely continuous across edges. Parameter
pairs outside these charts have a positive exact-image separation by
compactness; uniform \(C^0\) closeness preserves it. Consequently
the numerical map remains an embedding with nondegenerate elements.
This also justifies the normal graph used on the safe interval.

For fixed \(h\), the bounded nodal positions, positive element
Jacobian lower bounds, and \eqref{eq:cubic-actual-schur} keep the
trajectory in a compact subset of the ODE domain. The vector field
is bounded there, so a finite endpoint has a limiting nodal value
from which local existence restarts. Strict bootstrap improvement
excludes an endpoint before \(T\). Equation \eqref{eq:cubic-ODE}
also gives uniqueness of the velocity and multiplier.

Finally the original closest lift and
\eqref{eq:cubic-interpolation}, \eqref{eq:cubic-map-energy} give
\[
 \norm{X_h^\ell-X}_{L^2(\Gamma^0)}\leq CE+Ch^4
                      \leq C_T h^3\sqrt\ell.
\]
The signed-distance estimate follows from
\eqref{eq:cubic-bulk-distance}, \eqref{eq:cubic-bulk-improvement}.
All endpoint estimates hold by continuity. This proves
Theorem~\ref{thm:cubic}.

\begin{remark}\label{rem:quadratic-open}
These results partially solve the low-degree question: degree one
fails for an admissible family, and degree three converges for every
original admissible family under the smooth-solution hypothesis.
Degree two remains unresolved. Even substituting a putative
\(F=O(h^2)\) in \eqref{eq:cubic-tangent-energy} gives only the driving
bound \(O(h^2\ell)\), hence a map bound \(O(h\sqrt\ell)\).
The inverse estimate then cannot exclude local element degeneration.
This failure of the present estimate is not a quadratic counterexample.
\end{remark}

\begin{thebibliography}{9}
\bibitem{GLT2026}
G.~Gao, B.~Li and R.~Tang,
\emph{Convergent finite element approximations of surface evolution
with relaxed minimal deformation},
Numerische Mathematik \textbf{158} (2026), 671--714.
\href{https://doi.org/10.1007/s00211-026-01529-3}{doi:10.1007/s00211-026-01529-3}.
Author manuscript:
\url{https://libuyang.com/NM-15.pdf}.
\bibitem{TaoFourier}
T.~Tao, \emph{245C, Notes 2: The Fourier transform},
lecture notes, 6 April 2009, Theorem~43.
\url{https://terrytao.wordpress.com/2009/04/06/the-fourier-transform/}.
\end{thebibliography}

\section*{AI-use statement}
This work was developed with substantial assistance from OpenAI Codex
in mathematical exploration, algebraic checking and manuscript preparation.

\end{document}
